%!TEX TS-program = xelatex
%!TEX root = ../KTheorie_SS15/K-Theorie.tex
\RequirePackage[thinlines]{easybmat} %-- muss aufgrund eines Bugs vor etex und tikz geladen werden
\documentclass[a4paper, twoside, headsepline, index=totoc,toc=listof, fontsize=10pt, cleardoublepage=empty, headinclude, DIV=12, BCOR=5mm, titlepage,draft]{scrartcl}
\usepackage{scrtime} % KOMA, Uhrzeit ermoeglicht
\usepackage{scrpage2} % wie fancyhdr, nur optimiert auf KOMA-Skript, leicht andere Syntax
\usepackage{etoolbox}
\usepackage{letltxmacro}
\usepackage{ifthen}

%--Farbdefinitionen
%-- muss vor tikz geladen werden
\usepackage[usenames, table, x11names]{xcolor}
\definecolor{dark_gray}{gray}{0.45}
\definecolor{light_gray}{gray}{0.6}
\usepackage[final]{graphicx}

%--Zum Zeichnen
%-- muss vor polyglossia bzw. babel geladen werden
\usepackage{tikz}
\usepackage{tikz-cd}
\usetikzlibrary{external}
\tikzset{>=latex}
\usetikzlibrary{shapes,arrows,intersections}
\usetikzlibrary{calc,3d}
\usetikzlibrary{decorations.pathreplacing,decorations.markings, decorations.pathmorphing}
\usetikzlibrary{angles}

%-- Konfiguration von tikzexternalize
\tikzexternalize[prefix=tikz/,up to date check=diff]
\pgfkeys{/pgf/images/include external/.code=\includegraphics{#1}}
\tikzset{external/system call={lualatex \tikzexternalcheckshellescape -halt-on-error -interaction=batchmode --shell-escape -jobname "\image" "\texsource"}}


%-- tikzexternalize fuer tikzcd deaktivieren, da inkompatibel
\AtBeginEnvironment{tikzcd}{\tikzexternaldisable}
\AtEndEnvironment{tikzcd}{\tikzexternalenable}

%-- um Inkompatibilitaeten von quotes und polyglossia bzw. babel zu vermeiden
\tikzset{
  every picture/.append style={
    execute at begin picture={\shorthandoff{"}},
    execute at end picture={\shorthandon{"}}
  }
}
\usetikzlibrary{quotes}
\usepackage{pgfplots}
\usepgfplotslibrary{colormaps}
\pgfplotsset{compat=1.12}



%-- Mathesymbole etc.
\usepackage{mathtools} % beinhaltet amsmath
\mathtoolsset{showonlyrefs, centercolon}
\newtagform{brackets}[\textbf]{[}{]}
\usetagform{brackets}
\usepackage{wasysym} % zusätzliche Symbole
\usepackage{amssymb} % zusätzliche Symbole
% \usepackage{latexsym} % zusätzliche Symbole
\usepackage{stmaryrd} % für Blitz
\usepackage{nicefrac} % schräge Brüche
\usepackage{faktor}
\newcommand{\Faktor}[1]{\faktor[\textstyle]{#1}}
\usepackage{cancel} % Befehle zum Durchstreichen
\usepackage{mathdots} % Verbesserung von Punkten wie zB \ldots
\usepackage{mathrsfs} % das X wurde benötigt

%-- einzelne Symbole aus dem mathx-package
\DeclareFontFamily{U}{mathx}{\hyphenchar\font45}
\DeclareFontShape{U}{mathx}{m}{n}{<-> mathx10}{}
\DeclareSymbolFont{mathx}{U}{mathx}{m}{n}
\DeclareMathSymbol{\bigplus}{1}{mathx}{"90}

%-- 
\DeclareSymbolFont{bbold}{U}{bbold}{m}{n}
\DeclareSymbolFontAlphabet{\mathbbold}{bbold}
\newcommand{\mathds}[1]{\mathbb{#1}} % Um Kompatibilitaet mit frueheren Benutzung von dsfont herzustellen bzw zusätzliche Kapselung
\newcommand{\ind}{\mathbbold{1}} % charakteristische-Funktion-Eins

\def\mathul#1#2{\color{#1}\underline{{\color{black}#2}}\color{black}} %farbiges Untersteichen im Mathe-Modus
\renewcommand{\le}{\leqslant}
\renewcommand{\ge}{\geqslant}

%-- Underbrace als Befehl in LaTeX-Syntax (und ohne Spacingprobleme mit nachfolgenden Operatoren...)
\newcommand{\Underbrace}[2]{{\underbrace{#1}_{#2}}}

%-- Alles was mit Schrift und XeTeX zu tun hat
\usepackage[euler-digits]{eulervm}
\usepackage[no-math]{fontspec}
\usepackage{polyglossia} % moderner babel-ersatz
\setmainlanguage[spelling=new,babelshorthands=true]{german}
\shorthandoff{"}
\setotherlanguage{english}
% \newcommand\glqq{"}
% \newcommand\grqq{"}
\defaultfontfeatures{Mapping=tex-text, WordSpace={1.2}, Ligatures={Required,Common,Contextual}} %

%-- Alle hier aufgeführten Schriftarten sind in jeder TeXlive-Installation vorhanden!!
%-- Hauptschriftart festlegen
\setmainfont{texgyrepagella}[Extension=.otf,UprightFont=*-regular,BoldFont=*-bold,ItalicFont=*-italic,BoldItalicFont=*-bolditalic,ItalicFeatures={Style=Historic}]
% \setmainfont[BoldFont={* Bold}, ItalicFont={* Light Italic},SmallCapsFont={Linux Libertine O}, SmallCapsFeatures={Letters=SmallCaps}]{Source Sans Pro}


%-- Serifenlose Schriftart festlegen
% \setsansfont{Roboto}[Scale=MatchUppercase,Extension=.ttf,UprightFont=*-Regular,BoldFont=*-Medium,ItalicFont=*-RegularItalic,BoldItalicFont=*-MediumItalic]
% \setsansfont{RobotoSlab}[Scale=MatchUppercase,Extension=.ttf,UprightFont=*-Regular,BoldFont=*-Bold,ItalicFont=*-Thin,SmallCapsFont=*-Regular]
% \setsansfont{FiraSans}[Extension=.otf, UprightFont=*-Regular, BoldFont=*-Medium, ItalicFont=*-RegularItalic, BoldItalicFont=*-MediumItalic]
% \setsansfont{SourceSansPro}[Scale=MatchUppercase,Extension=.otf,UprightFont=*-Regular,BoldFont=*-Semibold,ItalicFont=*-LightIt,BoldItalicFont=*-SemiboldIt]
% \setsansfont{texgyreadventor}[Scale=MatchUppercase, Extension=.otf, UprightFont=*-regular, BoldFont=*-bold, ItalicFont=*-italic,BoldItalicFont=*-bolditalic]
% \setsansfont{Kurier}[Scale=MatchUppercase, Extension=.otf, UprightFont=*Medium-Regular, BoldFont=*Heavy-Regular, ItalicFont=*Medium-Italic,BoldItalicFont=*Heavy-Italic]
% \setsansfont{Iwona}[Scale=MatchUppercase, Extension=.otf, UprightFont=*Medium-Regular, BoldFont=*Heavy-Regular, ItalicFont=*Medium-Italic,BoldItalicFont=*Heavy-Italic]
\setsansfont{LinBiolinum}[Scale=MatchUppercase, Extension=.otf, UprightFont=*_R, BoldFont=*_RB, ItalicFont=*_RI,BoldItalicFont=*_RBO]


% \newfontfamily{\geometric}{texgyreheros}[Scale=MatchLowercase, Extension=.otf, UprightFont=*-regular, BoldFont=*-bold, ItalicFont=*-italic]
% \addtokomafont{disposition}{\geometric}
\setmonofont{SourceCodePro}[Scale=0.9,Extension=.otf,UprightFont=*-Regular, BoldFont=*-Semibold, ItalicFont=*-Light]
\usepackage{xltxtra}
\usepackage{fontawesome}
\usepackage[final]{microtype}


%--Mixed
\usepackage[shortlabels,inline]{enumitem}
\setlist[itemize,1]{label=\faCaretRight}
\setlist[enumerate]{font=\bfseries}
\usepackage[german=quotes]{csquotes}
\usepackage{makeidx}
\usepackage{booktabs}
\usepackage{wrapfig}
\usepackage{float}
\usepackage[margin=10pt, font=small, labelfont={sf, bf}, format=plain, indention=1em]{caption}
\captionsetup[wrapfigure]{name=Abb. }
\usepackage{stackrel}
\usepackage{multicol}

\flushbottom


%--Unterstreichung
\usepackage[normalem]{ulem}
\setlength{\ULdepth}{1.8pt}

%--Indexverarbeitung
\newcommand{\bet}[1]{\emph{#1}}
\newcommand{\Index}[1]{\bet{#1}\index{#1}}
\makeindex
\setindexpreamble{{\noindent \itshape Die \emph{Seitenzahlen} sind mit Hyperlinks zu den entsprechenden Seiten versehen, also anklickbar} \faHandUp \par \bigskip}
\renewcommand{\indexpagestyle}{scrheadings}


%--Marginnote & todonotes
\deffootnote[1.5em]{1.5em}{1.5em}{\textsuperscript{\thefootnotemark}\ }
\usepackage[fulladjust]{marginnote}
\renewcommand*{\marginfont}{\color{dark_gray} \itshape \footnotesize}
\usepackage[textsize=small]{todonotes}
\usepackage{ragged2e}
\renewcommand*{\raggedleftmarginnote}{\RaggedLeft}
\renewcommand*{\raggedrightmarginnote}{\RaggedRight}

%--Todonotes mit tikz/externalize kompatibel machen
\LetLtxMacro{\oldtodo}{\todo}
\renewcommand{\todo}[2][]{\tikzexternaldisable\oldtodo[#1]{#2}\tikzexternalenable}
\LetLtxMacro{\oldmissingfigure}{\missingfigure}
\renewcommand{\missingfigure}[2][]{\tikzexternaldisable\oldmissingfigure[{#1}]{#2}\tikzexternalenable}

%--Konfiguration von Hyperref pdfstartview=FitH, 
\usepackage[hidelinks, pdfpagelabels,  bookmarksopen=true, bookmarksnumbered=true, linkcolor=black, urlcolor=SkyBlue2, plainpages=false,pagebackref, citecolor=black, hypertexnames=true, pdfauthor={Jannes Bantje}, pdfborderstyle={/S/U}, linkbordercolor=SkyBlue2, colorlinks=false,final]{hyperref}


\newcommand{\appendLink}[1]{#1\,\faExternalLink}
\newcommand{\hrefsym}[2]{\href{#1}{\texttt{\appendLink{#2}}}}
\newcommand{\hrefsymmail}[2]{\href{#1}{\texttt{\faEnvelope\,#2}}}
\renewcommand{\url}[1]{\hrefsym{#1}{\nolinkurl{#1}}}



% -- QR-Codes
\usepackage{qrcode} %-- hinter hyperref laden!

%--Römische Zahlen
\newcommand{\RM}[1]{\MakeUppercase{\romannumeral #1{}}}

%%--Abkürzungen etc.
\newcommand{\light}{\text{\Large $\lightning$}}

%-- Definitionen von weiteren Mathe-Befehlen
\DeclareMathOperator{\re}{Re} %Realteil
\DeclareMathOperator{\im}{Im} %Imaginaerteil
\DeclareMathOperator{\id}{id} %identische Abbildung
\DeclareMathOperator{\Sp}{Sp} %Spur
\DeclareMathOperator{\sgn}{sgn} %Signum
\DeclareMathOperator{\alt}{Alt} %Alternierende n-linearFormen
\DeclareMathOperator{\End}{End} %Endomorphismen
\DeclareMathOperator{\Vol}{Vol} %Volumen
\DeclareMathOperator{\dom}{dom} %Domain
\DeclareMathOperator{\Hom}{Hom} %Homomorphismen
\DeclareMathOperator{\Iso}{Iso} %Isomorphismen
\DeclareMathOperator{\END}{End} %Isomorphismen
\DeclareMathOperator{\bild}{Bild} %Bild
\DeclareMathOperator{\Kern}{Kern}
\DeclareMathOperator{\ev}{ev} %Auswertungsabbildung
\DeclareMathOperator{\rg}{rg} %Rang
\DeclareMathOperator{\Rg}{Rg} %Rang
\DeclareMathOperator{\diam}{diam} %Durchmesser
\DeclareMathOperator{\dist}{dist} %Distanz
\DeclareMathOperator{\grad}{grad} %Gradient
\DeclareMathOperator{\dive}{div} %Gradient
\DeclareMathOperator{\rot}{rot} %Rotation
\DeclareMathOperator{\hess}{Hess} %Hesse-Matrix
\DeclareMathOperator{\koker}{Koker} %Kokern
\DeclareMathOperator{\coker}{coker} %Kokern
\DeclareMathOperator{\aut}{Aut} %Automorphismen
\DeclareMathOperator{\ord}{ord} %Ordnung
\DeclareMathOperator{\ggT}{ggT} %ggT
\DeclareMathOperator{\kgV}{kgV} %kgV
\DeclareMathOperator{\Gr}{Gr} %Gerade
\DeclareMathOperator{\Kr}{Kr} %Kreis
\DeclareMathOperator{\Char}{char} %Charakteristik
\DeclareMathOperator{\Aut}{Aut} %Automorphismen
% \DeclareMathOperator{\D}{D} %formale Ableitung
\newcommand{\D}{\ensuremath{\mathrm{D}\mkern-1.5mu}}
\newcommand{\mathd}{\ensuremath{\mathrm{d}\mkern-1.0mu}}
\newcommand{\Tmap}{\ensuremath{\mathrm{T}\mkern-0.85mu}}
\newcommand{\opL}{\ensuremath{\mathrm{L}\mkern-0.6mu}}
\DeclareMathOperator{\cov}{cov} %Kovarianz
\DeclareMathOperator{\Gal}{Gal} %Galois
\DeclareMathOperator{\supp}{supp} %Träger
\DeclareMathOperator{\Sym}{Sym} %Symmetrische Gruppe
\DeclareMathOperator{\Zyl}{Zyl} %Zylinder
\DeclareMathOperator{\Mod}{Mod} %Moduln
\DeclareMathOperator{\EPK}{EPK} %Einpunktkompaktifizierung
\DeclareMathOperator{\conj}{conj} %Konjugation
\DeclareMathOperator{\ann}{ann} %Annulator
\DeclareMathOperator{\tor}{tor} %Torsionsmodul
\DeclareMathOperator{\Int}{Int} %Interior/Inneres
\DeclareMathOperator{\Br}{Br} %Brauerergruppe
\DeclareMathOperator{\GL}{GL} %allgemeine lineare Gruppe
\DeclareMathOperator{\SL}{SL} %spezielle lineare Gruppe
\DeclareMathOperator{\SU}{SU} %spezielle unitäre Gruppe
\DeclareMathOperator{\Tr}{Tr} %Spur/Trace
\DeclareMathOperator{\conv}{conv} %Konvexe Teilmenge
\DeclareMathOperator{\Span}{span} %span
\DeclareMathOperator{\Diff}{Diff} %Menge der Diffeomorphismen
\DeclareMathOperator{\Isom}{Isom} %Menge der Isometrien
\DeclareMathOperator{\spec}{spec} %Spektrum
\DeclareMathOperator{\res}{res} %Resolventenmenge
\DeclareMathOperator{\resolv}{R} %Resolventenmenge
\DeclareMathOperator{\idx}{ind} %Fredholm Index
\DeclareMathOperator{\Fred}{Fred} %menge der Fredholm-Operatoren

\newcommand{\tr}{\mathrm{tr}}
\newcommand{\w}{\mathrm{w}}
\newcommand{\sa}{\mathrm{s.a.}}
\newcommand{\so}{\mathrm{s.o.}}
\newcommand{\weakT}[1]{\ensuremath{\mathcal{T}_{#1}^{\mkern+1.0mu\text{\raisebox{0.4ex}{$\mathrm{w}$}}}}}
\newcommand{\weakTstar}[1]{\ensuremath{\mathcal{T}_{#1}^{\mkern+1.0mu\text{\raisebox{0.4ex}{$\mathrm{w}$}}^*}}}
\newcommand{\finSub}{\subset\mkern-0.7mu \subset}
\newcommand{\rand}[1]{\ensuremath{\partial^{\scriptscriptstyle #1}}}


% -- Kategorien
\DeclareMathOperator{\Ob}{Ob} %Objekte
\DeclareMathOperator{\obj}{obj} %Objekte
\DeclareMathOperator{\Mor}{Mor} %Morphismen
\newcommand{\stiso}{\mathrm{st. iso}} % stabiler Isomorphismus
\newcommand{\st}{\mathrm{st}}

\newcommand{\TOP}{\textsc{Top}}
\newcommand{\HTOP}{\textsc{HTop}}
\newcommand{\VR}{\textsc{VR}}
\newcommand{\MOD}{\textsc{Mod}}
\newcommand{\MONOIDE}{\textsc{Monoide}}
\newcommand{\MENGEN}{\textsc{Mengen}}
\newcommand{\MAN}{\textsc{Man}}
\newcommand{\GRUPPEN}{\textsc{Gruppen}}
\newcommand{\ABELGRUPPEN}{\textsc{Abel.Gruppen}}
\newcommand{\KAT}{\textsc{Kat}}
\newcommand{\FUN}{\textsc{Fun}}
\newcommand{\SIMP}{\textsc{Simp}}
\newcommand{\VEKT}{\textsc{Vekt}}


\newcommand{\op}{\mathrm{op}}
\newcommand{\ab}{\mathrm{ab}}
\newcommand{\pr}{\mathrm{pr}}
\newcommand{\pt}{\mathrm{pt}}
\newcommand{\const}{\mathrm{const}}
\newcommand{\CW}{\ensuremath{\mathrm{CW}}}
\DeclareMathOperator{\cell}{cell}
\DeclarePairedDelimiter{\homologieklasse}{\llbracket}{\rrbracket}

\newcommand{\Ker}{\ensuremath{\Kern \,}}
\newcommand{\zirkel}{\rotatebox[origin = c]{-90}{$\varangle$}}
\newcommand{\wirkung}{\!\curvearrowright\!}

\newcommand{\oE}{\OE}


%--Beweisende
\newcommand{\bewende}{\ifmmode \tag*{$\square$} \else \hfill $\square$ \fi}


% -- theorem packages
\usepackage{amsthm}
\usepackage{thmtools,thm-restate}
\renewcommand{\listtheoremname}{Übersicht aller Aussagen}

% -- Theoreme als PDF-Lesezeichen
\usepackage{bookmark}
\bookmarksetup{open,numbered}
\makeatletter
\newcommand*{\theorembookmark}{%
  \bookmark[
    dest=\@currentHref,
    rellevel=1,
    keeplevel,
  ]{%
    \thmt@thmname\space\csname the\thmt@envname\endcsname
    \ifx\thmt@shortoptarg\@empty
    \else
      \space(\thmt@shortoptarg)%
    \fi
  }%
}   
\makeatother

% -- Definition der einzelnen Umgebungen
\declaretheoremstyle[%
	headfont=\sffamily\bfseries,
	notefont=\normalfont\sffamily\scshape,
	bodyfont=\normalfont,
	headformat=\NUMBER\ \NAME\NOTE,
	headpunct=.,
	postheadspace=1em,
	spaceabove=15pt,spacebelow=10pt,
	postheadhook=\theorembookmark]%
{mainstyle}
\declaretheoremstyle[%
	headfont=\sffamily\bfseries,
	notefont=\normalfont\sffamily\scshape,
	bodyfont=\normalfont,
	headformat=\NUMBER\NAME\NOTE,
	headpunct=.,
	postheadspace=1em,
	spaceabove=15pt,spacebelow=10pt,
	postheadhook=\theorembookmark]%
{mainstyle_unnumbered}
\declaretheoremstyle[%
	headfont=\sffamily\bfseries,
	notefont=\normalfont\sffamily\scshape,
	bodyfont=\normalfont,
	headformat=swapnumber,
	headpunct=.,
	postheadspace=1em,
	spaceabove=15pt,spacebelow=10pt,
	postheadhook=\theorembookmark,
	qed=\qedsymbol]%
{mainstyleB}
\declaretheorem[name=Definition,parent=section,style=mainstyle]{definition}
\declaretheorem[name=Definition,numbered=no,style=mainstyle_unnumbered]{definition*}
\declaretheorem[name=Theorem,sharenumber=definition,style=mainstyle]{theorem}
\declaretheorem[name=Theorem,numbered=no,style=mainstyle_unnumbered]{theorem*}
\declaretheorem[name=Proposition,sharenumber=definition,style=mainstyle]{proposition}
\declaretheorem[name=Lemma,sharenumber=definition,style=mainstyle]{lemma}
\declaretheorem[name=Satz,sharenumber=definition,style=mainstyle]{satz}
\declaretheorem[name=Korollar,sharenumber=definition,style=mainstyle]{korollar}
\declaretheorem[name=Korollar,sharenumber=definition,style=mainstyleB]{korollarB}


% -- Beweise
\declaretheoremstyle[headfont=\sffamily\itshape\scshape,bodyfont=\normalfont,headpunct=:,postheadspace=1em,spacebelow=12pt,spaceabove=2pt,qed=\qedsymbol]{beweise}
\declaretheoremstyle[headfont=\sffamily\bfseries,bodyfont=\normalfont,headpunct=:,postheadspace=1em,spacebelow=10pt,spaceabove=10pt]{bemerkungen}
\declaretheorem[name=Beweis,numbered=no,style=beweise]{beweis}
\declaretheorem[name=Bemerkung,numbered=no,style=bemerkungen]{bemerkung}
\declaretheorem[name=Beispiel,numbered=no,style=bemerkungen]{beispiel}
\declaretheorem[name=Übung,numbered=no,style=bemerkungen]{uebung}
\declaretheorem[name=Erinnerung,numbered=no,style=bemerkungen]{erinnerung}

%--Volumen
\newcommand{\vol}[1]{\ensuremath{\Vol_{#1}}}

%--Skalarprodukt
\DeclarePairedDelimiterX\sprod[2]{\langle}{\rangle}{#1\,\delimsize\vert\,#2}
\DeclarePairedDelimiterX\skal[2]{\langle}{\rangle}{#1\,,\,#2}

%--Betrag, Gaußklammer
\DeclarePairedDelimiter{\abs}{\lvert}{\rvert}
\DeclarePairedDelimiter{\floor}{\lfloor}{\rfloor}
\DeclarePairedDelimiter{\ceil}{\lceil}{\rceil}

%--Norm
\DeclarePairedDelimiter\doppelstrich{\Vert}{\Vert}
\newcommand{\norm}[2][\relax]{
\ifx#1\relax \ensuremath{\doppelstrich*{#2}}
\else \ensuremath{\doppelstrich*{#2}_{#1}}
\fi}


%--Umklammern
\DeclarePairedDelimiter\enbrace{(}{)}
\DeclarePairedDelimiter\benbrace{[}{]}
\DeclarePairedDelimiter\lenbrace{<}{>}
\newcommand{\ssbrace}[1]{{\scriptscriptstyle\enbrace{#1}}}

%--Mengen
\DeclarePairedDelimiterX\mengenA[1]{\lbrace}{\rbrace}{#1}
\DeclarePairedDelimiterX\mengenB[2]{\lbrace}{\rbrace}{#1\, \delimsize\vert \, #2}

\makeatletter
\newcommand{\set}[2][\relax]{
\ifx#1\relax \ensuremath{
\mengenA*{#2}}
\else \ensuremath{%
  \mengenB*{#1}{#2}}
\fi}
\makeatother

\DeclareRobustCommand{\minwidthbox}[2]{%
  \ifmmode
    \expandafter\mathmakebox
  \else
    \expandafter\makebox
  \fi
  [\ifdim#2<\width\width\else#2\fi]{#1}%
}


%--offen und abgeschlossen
\newcommand{\off}{\! \stackrel[\text{offen}]{}{\subset} \!}
\newcommand{\abg}{\! \stackrel[\text{abg.}]{}{\subset} \!}

%--Differential
\newcommand{\diff}[2]{\ensuremath{\frac{{\partial #1}}{{\partial #2}} }}
\newcommand{\diffd}[2]{\ensuremath{\frac{\mathd #1}{\mathd #2} }}

%--Inhaltsverzeichnis
\usepackage[tocindentauto]{tocstyle}
\usetocstyle{KOMAlike}	
\shorthandon{"}
		
%--Punkte (als letztes laden)
\usepackage{ellipsis}

\newcommand{\fach}{$K$-Theorie und die Hopf-Invariante}
\newcommand{\semester}{SoSe 2015}
\newcommand{\homepage}{https://wwwmath.uni-muenster.de/reine/u/topos/lehre/SS2015/KTheorie-Hopf/Hopf.html}

\newcommand{\prof}{Dr.\ Ulrich Penning}
\publishers{\scalebox{14}{$K$}}
\providecommand{\verfasser}{Jannes Bantje}

%--Konfiguration von scrheadings
\setheadsepline{1pt}[\color{light_gray}]
\pagestyle{scrheadings}
%\fancyhf[H, F]{}
\clearscrheadfoot

\providecommand{\shortFach}{\fach}
\lehead{ \includegraphics[height=0.5 cm,keepaspectratio]{../!config/Bilder/Logo_WWU_Muenster_light_gray.pdf}}
\rehead{\rule{0cm}{0.5cm}\footnotesize \sffamily \color{light_gray} \verfasser -- Mitschrift \shortFach}
\lohead{\rule{0cm}{0.5cm} \footnotesize \sffamily \color{light_gray} Stand: \today \; \thistime[:]}
\rohead{\includegraphics[height=0.5 cm,keepaspectratio]{../!config/Bilder/fb10logo_gray.pdf}}


\ofoot[{ \color{dark_gray} \LARGE \sffamily \thepage}]{{ \color{dark_gray} \LARGE \sffamily \thepage}} %hier wir auch der plain Stil bearbeitet!
\automark{section}
\ifoot{ \color{dark_gray} \small \leftmark}

%--Metadaten
\author{\verfasser}
\titlehead{\includegraphics[height=1.5cm, keepaspectratio]{../!config/Bilder/Logo_WWU_Muenster.pdf}%
\hfill \includegraphics[height=1.3cm, keepaspectratio]{../!config/Bilder/fb10logo.pdf}}
\title{Skript \fach}
\subtitle{Mitschrift der Vorlesung  \enquote{\fach} von \prof}

\publishers{
Aktuelle Version verfügbar bei: \bigskip\\
\normalsize
\begin{minipage}[t]{0.45\textwidth}
	\centering{
	\qrcode[height=3.5cm, version=4]{https://github.com/JaMeZ-B/latex-wwu}
	\medskip\\
	\includegraphics[height=0.4cm, keepaspectratio]{../!config/Bilder/github_octo.pdf} 
	\includegraphics[height=0.4cm, keepaspectratio]{../!config/Bilder/GitHub_Logo.pdf} (inklusive Sourcecode)\\
	{\small\url{https://github.com/JaMeZ-B/latex-wwu}}}
\end{minipage}
\hfill
\begin{minipage}[t]{0.45\textwidth}
	\centering{
	\qrcode[height=3.5cm, version=4]{btsync://B6WH2DISQ5QVYIRYIEZSF4ZR2IDVKPN3I?n=latex_share}
	\medskip\\
	\includegraphics[height=0.4cm, keepaspectratio]{../!config/Bilder/bt_sync_logo.pdf} 
	{\large \textbf{Bittorrent} Sync} \\
	\texttt{B6WH2DISQ5QVYIRYIEZSF4ZR2IDVKPN3I}
	}
\end{minipage}
}


\begin{document}
\hyphenation{Di-vi-sions-al-ge-bra}
\pagenumbering{Roman}
\maketitle
\begin{abstract}
\section*{Vorwort --- Mitarbeit am Skript}
Dieses Dokument ist eine Mitschrift aus der Vorlesung \enquote{\fach, \semester}, gelesen von \prof. Der Inhalt entspricht weitestgehend dem Tafelanschrieb. Für die
Korrektheit des Inhalts übernehme ich keinerlei Garantie! Für Bemerkungen und Korrekturen -- und seien es nur Rechtschreibfehler -- bin ich sehr dankbar. 
Korrekturen lassen sich prinzipiell auf drei Wegen einreichen: 
\begin{itemize}
	\item Persönliches Ansprechen in der Uni, Mails an \hrefsymmail{mailto:\mail}{\mail} (gerne auch mit annotieren PDFs) 
	\item \emph{Direktes} Mitarbeiten am Skript: Den Quellcode poste ich auf GitHub (siehe oben), also stehen vielfältige Möglichkeiten der Zusammenarbeit zur Verfügung:
	Zum Beispiel durch Kommentare am Code über die Website und die Kombination Fork + Pull Request. Wer sich verdient macht oder ein Skript zu einer Vorlesung, die 
	ich nicht besuche, beisteuern will, dem gewähre ich gerne auch Schreibzugriff.
	
	Beachten sollte man dabei, dass dazu ein Account bei \url{github.com} notwendig ist, der allerdings ohne Angabe von persönlichen Daten angelegt werden kann. 
	Wer bei GitHub (bzw. dem zugrunde liegenden Open-Source-Programm \enquote{\texttt{git}}) -- verständlicherweise -- Hilfe beim Einstieg braucht, dem helfe ich gerne 
	weiter. Es gibt aber auch zahlreiche empfehlenswerte Tutorials im Internet.\footnote{zB. \url{https://try.github.io/levels/1/challenges/1}, ist auf Englisch, aber dafür 
	interaktives LearningByDoing}
	\item \emph{Indirektes} Mitarbeiten: \TeX-Dateien per Mail verschicken. 
	
	Dies ist nur dann sinnvoll, wenn man einen ganzen Abschnitt ändern möchte (zB. einen alternativen Beweis geben), da ich die Änderungen dann per Hand einbauen muss! Ich freue mich aber auch über solche Beiträge!
\end{itemize}
\section*{Vorlesungshomepage}
{\centering 
\begin{minipage}[c][][c]{\textwidth}
	\centering \qrcode[height=3cm]{\homepage} \medskip\\
	\footnotesize \url{\homepage}
\end{minipage}
\par}

\end{abstract}

\tableofcontents
\cleardoubleoddemptypage

\pagenumbering{arabic}
\setcounter{page}{1}

% \setcounter{section}{-1}
\section*{Einführung} % (fold)
\label{sec:0}

\begin{definition*}[{name=[Definition: Reelle Divisionsalgebra]},label=sub:01,numbered=no]
Eine $\mathds{R}$-Algebra $\mathcal{A}$ mit $1$, die nicht notwendigerweise assoziativ ist, heißt \Index{reelle Divisionsalgebra}, falls jedes Element $a\not= 0$ invertierbar 
ist.	
\end{definition*}

\begin{beispiel}
Die reellen Zahlen $\mathds{R}$, die komplexen Zahlen $\mathds{C}$, die Quaternionen\footnote{siehe auch \url{https://de.wikipedia.org/wiki/Quaternion}} $\mathds{H}$ und die sogenannten Cayley-Zahlen\footnote{auch Oktonionen oder reelle Okataven; siehe  \url{http://de.wikipedia.org/wiki/Oktave_(Mathematik)}} $\mathds{O}$. Dabei ist
$\mathds{O} \cong \mathds{H} \oplus \mathds{H}$, wobei die Multiplikation wie folgt definiert ist:
\[
	(a,b) \cdot (d,c) = (ac - d^* b, da + bc^*) 
\]
Dabei ist $d= x + iy + jz + kw$ und $d^* = x -iy -jz - kw$.
\end{beispiel}

\noindent Wir wollen in dieser Vorlesung das folgende Theorem über Divisionsalgebren beweisen:
\begin{theorem*}[name={Adams},restate=Adams]
Folgende Aussagen sind äquivalent:
\begin{enumerate}[(i)]
	\item $\mathds{R}^n$ besitzt die Struktur einer reellen Divisionsalgebra.
	\item Entweder ist $n=1$ oder $n \ge 2$ ist gerade und es gibt eine stetige Abbildung $f \colon S^{2n-1} \to S^n$ mit \Index{Hopf-Invariante} $1$.
	\item Es gilt $n \in \set{1,2,4,8}$.
\end{enumerate}
\end{theorem*}

% subsection 11 (end)

\subsection*{Geplanter Verlauf der Vorlesung} % (fold)
Wir werden diesen Satz mit Hilfe der sogenannten $K$-Theorie beweisen. Dazu müssen wir aber einiges an Vorarbeit leisten. Der geplante Verlauf ist deswegen wie folgt:
\begin{enumerate}[1.]
	\item Verallgemeinerte Kohomologietheorien
	\item Vektorbündel
	\item $K$-Theorie
	\begin{enumerate}[label*=\arabic*.]
		\item Produkte in $K$-Theorie
		\item Bott-Periodizität
		\item Thom-Isomorphismus
	\end{enumerate}
	\item Hopf-Invariante und der Beweis des Satzes
\end{enumerate}
% subsection inhalt (end)
% section 0 (end)
\newpage

\section{Kohomologietheorien} % (fold)
\label{sec:1}

\begin{definition}[{name=[Kontravarianter Funktor]}]
Seien $\mathcal{C}$ und $\mathcal{D}$ Kategorien. Ein \Index{kontravarianter Funktor} $F \colon \mathcal{C} \to \mathcal{D}$ 
\begin{itemize}
	\item ordnet jedem Objekt $c \in \mathrm{Obj}(\mathcal{C})$ ein Objekt $F(c) \in \mathrm{Obj}(\mathcal{D})$ zu,
	\item ordnet jedem Morphismus $f \colon c \to c'$ in $\mathrm{Mor}_{\mathcal{C}}(c,c')$ einen Morphismus $F(f) \colon F(c') \to F(c)$ zu,
\end{itemize}
sodass folgende Eigenschaften gelten:
\begin{itemize}
	\item Für $f \colon c \to c'$ und $g \colon c' \to c''$ gilt $F(g \circ f) = F(f) \circ F(g)$
	\item $F(\id_c) = \id_{F(c)}$
\end{itemize}
\end{definition}

\begin{beispiel}[name=Dualraumfunktor]
Sei $K$ ein Körper. Sei $\VEKT_K$ die Kategorie der endlich dimensionalen $K$-Vektorräume und linearen Abbildungen. Sei 
${ }^* \colon \VEKT_K \to \VEKT_K$ gegeben durch $V \mapsto \Hom_K(V,K)=V^*$. Für $f \colon V \to W$ sei $f^* \colon W^* \to V^*$, 
$\varphi \mapsto \varphi \circ f$. Dies ist ein kontravarianter Funktor. Insbesondere ist $(f \circ g)^* = g^* \circ f^*$.	
\end{beispiel}


\begin{definition}[{name=[Verallgemeinerte Kohomologietheorie]}]
Sei $\TOP^2$ die Kategorie der Raumpaare $(X,A)$. Sei $V \colon \TOP^2 \to \TOP^2$ der Funktor $(X,A) \mapsto(A,\emptyset)$. Sei $R$ ein kommutativer Ring mit $1$. Eine 
\Index{verallgemeinerte Kohomologietheorie} $h^* = (h^n, \partial^n)_{n \in \mathds{Z}}$ mit Werten in $R$-Moduln ist ein kontravarianter Funktor 
$h^* \colon \TOP^2 \to \textsc{Gr}\text{-}R\text{-}\MOD$ zusammen mit einer natürlichen Transformation 
\[
	\partial^* \colon h^* \circ V \to h^{*+1},
\]
so dass die folgenden Eigenschaften gelten:
\begin{enumerate}[a)]
	\item \Index{Homotopieinvarianz}: Es seien $f,g \colon (X,A) \to (Y,B)$ homotope Abbildungen von Raumpaaren. Dann gilt für alle $n \in \mathds{Z}$: 
	\[
		h^n(f)= h^n(g)
	\]
	\item \Index{Lange exakte Paarsequenz}: Sei $(X,A)$ ein Raumpaar. Seien $i \colon (A,\emptyset) \hookrightarrow (X,\emptyset)$ und 
	$j \colon (X,\emptyset) \hookrightarrow (X,A)$ die kanonischen Inklusionen. Dann ist die Folge 
	\[
		\begin{tikzcd}
			\cdots \rar["\partial^{n-1}"] & h^n(X,A;R) \rar["h^{n}(j)"] & h^n (X,\emptyset;R) \rar["h^n(i)"] & h^n(A, \emptyset;R) \rar["\partial^n"] & h^{n+1}(X,A;R) \rar & 
			\cdots
		\end{tikzcd}
	\]
	exakt.
	\item \Index{Ausschneidung}: Sei $(X,A)$ ein Raumpaar, $U \subset A$, sodass $\overline{U} \subseteq \mathring{A}$. Dann ist die von der Inklusion 
	$\iota \colon (X \setminus U, A \setminus U) \hookrightarrow (X,A)$ induzierte Abbildung $h^n(\iota) \colon h^n(X,A;R) \to h^n(X \setminus U, A \setminus U;R)$ ein 
	Isomorphismus.
\end{enumerate}	
\end{definition}

\begin{bemerkung}
\mbox{ }
\begin{enumerate}[(i)]
	\item Wir schreiben kurz $h^n(X,A)$ für $h^n(X,A;\mathds{Z})$ und $f^*$ für $h^n(f)$, falls sich der Grad aus dem Kontext ergibt. 
	\item $h^n(\set{\mathrm{pt}}) =: h^n$ heißen \Index{Koeffizienten der Kohomologietheorie}.
	\item Einige Kohomologietheorien besitzen ein \Index{externes Produkt}, das heißt eine natürliche Transformation
	\[
		h^n(X,A;R) \times h^m(Y,B;R) \longrightarrow h^{n+m} \enbrace[\big]{(X,A) \times (Y,B);R} 
	\]
	die bilinear und assoziativ ist. 
	Dabei ist $(X,A) \times (Y,B) = \enbrace[\big]{X \times Y, (A \times Y) \cup (X \times B)}$.
	
	Sei $\Delta \colon (X,A) \to (X,A) \times (X,A)$ die Diagonalabbildung.
	Falls ein externes Produkt für $h^*$ existiert, dann ist $h^*(X,A;R) = \bigoplus_{n \in \mathds{Z}} h^n(X,A;R)$ ein \Index{graduierter Ring} mit der Multiplikation 
	\[
		\begin{tikzcd}[row sep=large, column sep=small]
			h^n(X,A;R) \times h^m(X,A;R) \arrow[dashed]{rr} \arrow["\times"']{dr} & & h^{n+m}(X,A;R) \\
			& h^{n+m} \enbrace[\big]{(X,A) \times (X,A);R} \arrow["h^{n+m}(\Delta)"']{ur} &
		\end{tikzcd}
	\]
\end{enumerate}
\end{bemerkung}

\noindent Wir betrachten im Folgenden Kohomologietheorien, die auf einer Unterkategorie $\TOP^2_{\mathrm{Kof}}$ von $\TOP^2$ definiert sind:


\begin{definition}[{name=[Ko-Raumpaar und Kofaserung]},label=def_koraumpaar]
Ein Raumpaar $(X,A)$ heißt \Index{Ko-Raumpaar}, falls
\begin{itemize}
	\item $X$ ein kompakter Hausdorffraum ist,
	\item $A \subset X$ abgeschlossen ist,
	\item die Inklusion $A \hookrightarrow X$ eine \Index{Kofaserung} ist, das heißt für jede stetige Abbildung\marginnote{Homotopieerweiterungseigenschaft} 
	$f \colon A \times I \cup X \times \set{0} \to Y$ existiert 
	$H \colon X \times I \to Y$, sodass das folgende Diagramm kommutiert:
	\[
		\begin{tikzcd}[sep=large]
			A \times I \cup X \times \set{0} \rar["f"] \dar[hook]& Y \\
			X \times I \urar["H"'] & 
		\end{tikzcd}
	\] 
\end{itemize}
Mit $\TOP^2_{\mathrm{Kof}} \subseteq \TOP^2$ bezeichnen wir die Kategorie der Ko-Raumpaare und stetigen Abbildungen.	
\end{definition}

\begin{beispiel}
\leavevmode
\begin{itemize}
	\item Sei $X$ ein kompakter \CW-Komplex, $A \subset X$ ein Unterkomplex. Dann ist $\iota \colon A \hookrightarrow X$ eine Kofaserung.
	\item $X$ eine kompakte Mannigfaltigkeit, $A$ eine abgeschlossene Untermannigfaltigkeit.
	\item Falls $(X, \set{x_0})$ ein Objekt in $\TOP^2_{\mathrm{Kof}}$ ist, dann heißt $X$ \Index{wohlpunktiert}.
\end{itemize}	
\end{beispiel}

\noindent Zur Veranschaulichung des Begriffes \enquote{Kofaserung} betrachten wir den Raum
$X= [-1,1] \times [-1,1]$, den Punkt $x_0=(0,0)$ und setzen $Y=X$. Sei $\gamma \colon I \to X$ ein Pfad mit $\gamma(0)= (0,0)=x_0$. Wir erhalten mit $\gamma$ und $\id$ eine 
Abbildung $f \colon \set{x_0} \times I \cup  X \times \set{0} \to X$.
\[
	\begin{tikzpicture}[scale=1.5]
		\draw[DodgerBlue4,thick] (1,0,1) -- (1,0,-1) -- (-1,0,-1) -- (-1,0,1) -- cycle;
		\draw[very thick,fill] (0,0,0) circle[radius=0.04] node[below]{$x_0$} -- (0,0.6,0);
		
		\draw[-to,thick] (2,0) -- (4,0) node[midway,above]{$f$};
		
		\draw[DodgerBlue4,thick] (7,0,1) -- (7,0,-1) -- (5,0,-1) -- (5,0,1) -- cycle;
		\draw[very thick] (6,0,0) node[below]{$x_0$} .. controls +(-15:0.4) and +(185:0.2) .. (6.7,0,-0.8) node[midway,above,xshift=0.2]{$\gamma$};
		\filldraw[very thick] (6,0,0) circle[radius=0.04];
	\end{tikzpicture}
\]
Kofaserung heißt in diesem Fall, dass der Raum $X$ entlang des Pfades \enquote{mitgezogen} werden kann. Allgemeiner: $A \hookrightarrow X$ ist eine Kofaserung, falls sich 
jedes Paar aus einer Homotopie $h \colon A \times I \to Y$ und einer \enquote{Startabbildung} $X \times \set{0} \to Y$, die $h_0$ fortsetzt, zu einer Homotopie 
$H \colon X \times I \to Y$ erweitern lässt.
% section 1 (end)
\newpage

\section{Vektorbündel} % (fold)
\label{sec:2}
Wir betrachten nur reelle und komplexe Vektorbündel, es sei also $K \in \set{\mathds{R},\mathds{C}}$.

\begin{definition}[{name=[$K$-Vektorbündel]},label=def_vektbuendel]
Sei $X$ ein topologischer Raum. Eine stetige surjektive Abbildung $p \colon E \to X$ heißt $K$-\Index{Vektorbündel}, falls gilt: 
\begin{enumerate}[a)]
	\item $ E_x := p^{-1}(\set{x})$ ist ein $K$-Vektorraum für alle $x \in X$. $E_x$ heißt die \Index{Faser} von $E$ über $x$. 
	\item Für alle Punkte $x \in X$ gibt es eine Umgebung $U$, ein $n \in \mathds{N}$ und einen Homöomorphismus $\varphi_U \colon p^{-1}(U) \to U \times K^n$, sodass das 
	folgende 
	Diagramm kommutiert
	\[
		\begin{tikzcd}
			p^{-1}(U) \arrow["\varphi_U"]{rr} \drar["p"'] & & U \times K^n  \dlar["\mathrm{pr}_U"] \\
			& U 
		\end{tikzcd}
	\]
	und die Einschränkung $\varphi_U |_{E_x} \colon E_x \to K^n$ für alle $x \in U$ ein Isomorphismus von $K$"=Vektorräumen ist. Diese Eigenschaft wird oft als 
	\Index{lokale Trivialität} bezeichnet.
\end{enumerate}
\end{definition}

\begin{beispiel}
\leavevmode
\begin{itemize}
	\item Sei $X$ ein beliebiger topologischer Raum. Dann ist $X \times K^n \xrightarrow{\mathrm{pr}_X} X$ das \Index{triviale Vektorbündel} über $X$.
	Für jeden topologischen Raum existiert also ein Vektorbündel.
	\item Betrachte $S^n \subset \mathds{R}^{n+1}$. Sei $\Tmap S^n = \set[(x,v) \in S^n \times \mathds{R}^{n+1}]{\skal*{x}{v}=0}$. Definiere $p \colon \Tmap S^n \to S^n$ 
	durch $(x,v) \mapsto x$. Damit erhalten wir ein $\mathds{R}$-Vektorbündel über $S^n$, das \Index{Tangentialbündel}. 
	
	\emph{Übung: Zeige lokale Trivialität. Siehe Anhang \ref{sub:lokal_trivial_Sn}}
\end{itemize}	
\end{beispiel}

\begin{bemerkung}
Die Abbildung $x \mapsto \dim_K (E_x)$ ist stetig, also lokal konstant, das heißt über einem zusammenhängenden Raum hat ein Vektorbündel konstante Faserdimension. 	
\end{bemerkung}


\begin{definition}[{name=[Vektorraumbündelmorphismus]}]
Seien $p \colon E \to X$ und $p' \colon E' \to X'$ zwei $K$-Vektorbündel. Ein \bet{Morphismus von Vektorbündeln}\index{Morphismus!von Vektorbündeln} oder 
\Index{Vektorraumbündelmorphismus} von $E$ nach $E'$ ist ein Paar $(f, \overline{f})$ von stetigen Abbildungen $\overline{f} \colon E \to E'$ und $f \colon X \to X'$, sodass 
das folgende Diagramm kommutiert
\[
	\begin{tikzcd}[sep=large]
		E \rar["\overline{f}"] \dar["p"'] & E' \dar["p'"] \\
		X \rar["f"] & X'
	\end{tikzcd}
\]
und sodass die Einschränkung $\overline{f}|_{E_x} \colon E_x \to E'_{f(x)}$ für jedes $x \in X$ ein $K$-Vektorraumhomomorphismus ist. Wir bezeichnen die Kategorie der 
endlich dimensionalen $K$-Vektorraumbündel über einem topologischen Raum $X$ zusammen mit den Morphismen $(\id_X, \overline{f})$ mit $\VEKT_K(X)$.	
\end{definition}


\section*{Operationen mit Vektorbündeln}

\begin{definition}[{name=[Pullback eines Vektorbündels]}]
Seien $X$ und $Y$ topologische Räume. Sei $f \colon Y \to X$ eine stetige Abbildung und sei $p \colon E \to X$ ein Vektorbündel. Dann heißt
\[
	f^* E := \set[(y,v) \in Y \times E]{f(y)=p(v)\rule{0cm}{0.9em}} 
\]
zusammen mit der Abbildung $q \colon f^* E \to Y$, $(y,v) \mapsto y$ der \Index{Pullback}\marginnote{Der Pullback lässt sich auch allgemein auf kategorieller Ebene 
definieren} von $p \colon E \to X$ entlang von $f$. Mit $\overline{f}(y,v)=v$ kommutiert das 
folgende Diagramm:
\[
	\begin{tikzcd}[sep=large]
		f^* E \rar["\overline{f}"] \dar["q"'] & E \dar["p"] \\
		Y \rar["f"] & X
	\end{tikzcd}
\]	
\end{definition}
% subsection 23 (end)


\begin{lemma}[{name=[Universelle Eigenschaft des Pullbacks]}]
\leavevmode
\begin{enumerate}[a)]
	\item $q \colon f^* E \to Y$ ist ein Vektorbündel und $(\overline{f},f)$ ist ein Morphismus von Vektorbündeln. 
	\item $q \colon f^* E \to Y$ ist durch die folgende universelle Eigenschaft eindeutig (bis auf Isomorphie) charakterisiert:
	
	Sei $r \colon F \to Z$ ein Vektorbündel, $(\overline{g},g)$ ein Morphismus zwischen $r \colon F \to Z$ und $p \colon E \to X$. Sei weiter $g' \colon Z \to Y$ eine 
	stetige Abbildung mit $f \circ g' = g$. Dann existiert genau ein $G \colon F \to f^* E$, so dass $(G,g')$ ein Morphismus von Vektorbündeln ist und weiter 
	$(\overline{f},f) \circ (G,g')= (\overline{g},g)$ gilt.
	\[
		\begin{tikzcd}[sep=large]
			F \drar["\exists!\, G", dashed] \dar["r"'] \arrow[bend left,"\overline{g}"]{drr} & & \\
			Z \drar["g'"'] \arrow["g" near end]{drr} & f^* E \rar["\overline{f}"] \dar["q" pos=0.3,crossing over] & E \dar["p"]\\
			& Y \rar["f"'] & X
		\end{tikzcd}
	\]
\end{enumerate}	
\end{lemma}
\begin{beweis}
Sei $y\in Y$. Sei $U \subset X$ eine Umgebung von $f(y)$ in $X$, sodass eine lokale Trivialisierung $\varphi_U \colon p ^{-1}(U) \to U \times K^n$ existiert. Dann ist 
$V:= f^{-1}(U)$ eine Umgebung von $y$ und es gilt
\[
	q ^{-1}(V) = \set[(y,v) \in V \times p ^{-1}(U)]{f(y)=p(v)} 
\]
Wir definieren nun 
\begin{align*}
	\psi_V \colon \,\,q ^{-1}(V) &\longrightarrow V \times K^n &  (y,v)&\longmapsto \enbrace[\big]{y, \pr_{K^n} \circ \varphi_U(v)} \\
	\kappa_V \colon V \times K^n &\longrightarrow q ^{-1}(V) &  (y,w) &\longmapsto \enbrace*{y, \varphi_U^{-1} \enbrace*{f(y),w} } 
\end{align*}
Dann gilt $\psi_V \circ \kappa_V=\id_{V \times K^n}$ und $\kappa_V \circ \psi_V = \id_{q^{-1}(V)}$. Außerdem kommutiert das Diagramm aus \autoref{def_vektbuendel}. Also ist 
$\psi_V$ eine lokale Trivialisierung von $f^*E \to Y$ über $V$. \\
\emph{Der Beweis von b) ist eine (relativ einfache) Übungsaufgabe}\todo{eventuell hinzufügen}.
\end{beweis}

\noindent Sei $\mathcal{C}$ eine (topologische) Kategorie. Hieraus lässt sich eine neue Kategorie $\mathcal{C}^\op$ definieren:
\[
	\obj(\mathcal{C}^\op) := \obj(\mathcal{C}) \quad , \quad \Mor_{\mathcal{C}^\op}(c,d) := \Mor_{\mathcal{C}}(d,c)
\]
mit der neuen Komposition
\[
	\begin{tikzcd}[row sep=0.1em]
		\bullet \mbox{ :}\hspace{-4em} & \Mor_{\mathcal{C}^\op}(c',c'') \times \Mor_{\mathcal{C}^\op}(c,c') \rar & \Mor_{\mathcal{C}^\op}(c,c'') \\
		& (f,g) \rar[mapsto] & f \bullet g = g \circ f
	\end{tikzcd}
\]
(Kovariante) Funktoren $\mathcal{C}^\op \to \mathcal{D}$ entsprechen dann kontravarianten Funktoren $\mathcal{C}\to \mathcal{D}$. 

Sei $\VEKT_K = \VEKT_K(\pt)$ die Kategorie der endlichdimensionalen $K$-Vektorräume und linearen Abbildungen. Ein Funktor
\begin{equation}
	F \colon \Underbrace{\VEKT_K \times \ldots \times \VEKT_K}{r} \times \Underbrace{\VEKT_K^\op \times \ldots \times \VEKT_K^\op}{s} \longrightarrow \VEKT_K \label{eq:24}\tag{\#}
\end{equation}
heißt \bet{stetig}\index{stetiger Funktor}, falls die induzierte Abbildung
\begin{gather*}
	\Mor_{\VEKT_K}(V_1,V'_1) \times \ldots \times \Mor_{\VEKT_K}(V_r,V'_r) \times \Mor_{\VEKT_K^\op}(V_{r+1},V'_{r+1}) \times \ldots \times \Mor_{\VEKT_K^\op}(V_{r+s},V'_{r+s}) \\
	\longrightarrow \enspace\Mor_{\VEKT_K} \enbrace*{F(V_1, \ldots,V_r,V_{r+1}, \ldots , V_{r+s}), F (V'_1, \ldots , V'_r, V'_{r+1}, \ldots , V'_{r+s})} 
\end{gather*}
stetig ist. Dies klappt nur für Kategorien, für die eine Topologie auf den Morphismen existiert, was auf lineare Abbildungen zwischen endlichdimensionalen Vektorräumen aber 
natürlich zutrifft.
% subsection 24 (end)


\begin{lemma}[{name=[Induzierter Funktor eines stetigen Funktors]},label=ind_funk_vektorbuendel]
Sei $F$ ein stetiger Funktor wie in \eqref{eq:24} und sei $X$ ein topologischer Raum. Dann induziert $F$ einen Funktor
\[
	F_X \colon \Underbrace{\VEKT_K(X) \times \ldots \times \VEKT_K(X)}{r} \times \Underbrace{\VEKT_K(X)^\op \times \ldots \times \VEKT_K(X)^\op}{s} \longrightarrow \VEKT_K(X)
\]
der verträglich ist mit Pullbacks und $F_\pt=F$ erfüllt.
\end{lemma}
\begin{beweis}
Sei $m=r+s$. Definiere 
\(
	F_X \enbrace*{E^\ssbrace{1}, \ldots , E^\ssbrace{m}} = \coprod_{x \in X} F \enbrace*{E_x^\ssbrace{1}, \ldots , E_x^\ssbrace{m}}   
\)
als Menge. Sei $\set{U_i}_i$ eine offene Überdeckung von $X$, so dass lokale Trivialisierungen 
\[
	\varphi_i^\ssbrace{k} \colon (p^\ssbrace{k})^{-1}(U_i) \longrightarrow U_i \times K^{\ell_k}
\]
existieren. $F$ induziert nun -- da die $\varphi_i^\ssbrace{k}$ Homöomorphismen sind -- Bijektionen
\[
	F_{U_i} \enbrace*{\varphi_i^\ssbrace{1}, \ldots , \varphi_i^\ssbrace{m}} \colon F_{U_i} \enbrace*{(p^\ssbrace{1})^{-1}(U_i), \ldots , (p^\ssbrace{m})^{-1}(U_i)}
	\longrightarrow U_i \times F \enbrace*{K^{\ell_1}, \ldots , K^{\ell_m}}   
\]
Jetzt gibt es eine eindeutige Topologie auf $F_X \enbrace*{E^\ssbrace{1}, \ldots , E^\ssbrace{m}}$, so dass die eben definierten Abbildungen 
$\psi_i = F_{U_i} \enbrace*{\varphi_i^\ssbrace{1}, \ldots , \varphi_i^\ssbrace{m}}$ Homöomorphismen werden, wobei $U_i \times F \enbrace*{K^{\ell_1}, \ldots , K^{\ell_m}}$ 
die Produkttopologie trägt: Eine Menge $V \subset F_X \enbrace*{E^\ssbrace{1}, \ldots , E^\ssbrace{m}}$ ist offen, falls sie die Vereinigung von Mengen der Form 
$\psi_i^{-1}(V_i)$ mit $V_i \subset U_i \times F \enbrace*{K^{\ell_1}, \ldots , K^{\ell_m}}$ offen ist.

Damit diese Definition konsistent ist, müssen wir folgendes überprüfen: Sei $V \subset (U_i \cap U_j) \times F(K^{\ell_1}, \ldots , K^{\ell_m})$ offen. Wir haben zu zeigen, dass 
dann auch $\psi_j \circ \psi_i^{-1}(V)$ offen ist. Aber es gilt $\psi_j \circ \psi_i ^{-1}(x,v)= \enbrace[\big]{x, \psi_{ij}(x)(v)}$ für eine Abbildung
\[
	\psi_{ij} \colon U_i  \cap U_j \longrightarrow \End_K \enbrace*{F(K^{\ell_1}, \ldots , K^{\ell_m})} 
\]
Es gilt $\psi_{ij} = F \enbrace[\big]{\psi_{ij}^\ssbrace{1}, \ldots , \psi_{ij}^\ssbrace{m}}$ für \emph{stetige} Abbildungen 
$\psi_{ij}^\ssbrace{k} \colon U_i \cap U_j \to \End_K(K^{\ell_k})$. Da $F$ stetig ist, folgt, dass alle $\psi_{ij}$ stetig sind und folglich auch $\psi_j \circ \psi_i ^{-1}$. 
Somit ist $\psi_j \circ \psi_i ^{-1}(V)$ offen. Es ist damit klar, dass $F_X \enbrace*{E^\ssbrace{1}, \ldots , E^\ssbrace{m}} \to X$ ein Vektorbündel ist, denn wir haben lokale Trivialisierungen konstruiert. Ferner gilt per Definition für eine stetige Abbildung $f \colon Y \to X$
\[
	f^* F_X \enbrace*{E^\ssbrace{1}, \ldots , E^\ssbrace{m}} \cong F_Y \enbrace*{f^*E^\ssbrace{1}, \ldots , f^* E^\ssbrace{m}} \qedhere
\]	
\end{beweis}

\begin{beispiel}
\autoref{ind_funk_vektorbuendel} liefert uns viele bekannte Konstruktionen für Vektorbündel:
\begin{itemize}
	\item $\oplus \colon \VEKT_K \times \VEKT_K \to \VEKT_K$ ist stetig $\Rightarrow $ direkte Summe von Vektorbündeln\index{Vektorbündel!direkte Summe}. 
	\item $\otimes \colon \VEKT_K \times \VEKT_K \to \VEKT_K$ ist stetig $\Rightarrow $ Tensorprodukt von Vektorbündeln\index{Vektorbündel!Tensorprodukt}. 
	\item Dualisieren von Vektorräumen ${ }^* \colon \VEKT_K^\op \to \VEKT_K$ ist stetig $\Rightarrow $ duales Vektorbündel
	\item Der Funktor $\Hom \colon \VEKT_K \times \VEKT_K^\op \to \VEKT_K$, $(V,W) \mapsto \Hom_K(W,V)$ ist stetig $\Rightarrow$ 
	Homomorphismenbündel\index{Homomorphismenbündel}.
\end{itemize}	
\end{beispiel}

\noindent Seien $p \colon E \to X$ und $q \colon F \to X$ zwei $K$-Vektorbündel. Dann ist $E \times F \xrightarrow{(p,q)} X \times X$ auch ein Vektorbündel, wenn 
$E_x\times F_x$ die Vektorraumstruktur der äußeren direkten Summe trägt.

\begin{uebung}
Sei $\Delta  \colon X \to X \times X$ die Diagonalabbildung. Zeige, dass $\Delta^*(E \times F) \cong E \oplus F$ als Vektorbündel über $X$.	
\end{uebung}



\section*{Homotopieinvarianz von Pullbacks}
\begin{erinnerung}
Ein topologischer Raum heißt \bet{normal}\index{normaler topologischer Raum}, falls sich zwei disjunkte abgeschlossene Mengen durch offene Mengen trennen lassen. 	
Ein wichtiger Satz über normale Räume ist der folgende:
\end{erinnerung}

\begin{satz}[{name={Fortsetzungssatz von Tietze}},label=fort_tietze]
Sei $X$ ein normaler Raum, $A \subset X$ abgeschlossen und $V$ ein endlich dimensionaler $K$-Vektorraum. Zu jeder stetigen Abbildung $f \colon A \to V$ existiert eine 
stetige Abbildung $F \colon X \to V$ mit $F|_A=f$.
\end{satz}


\begin{bemerkung}
Jeder kompakte Hausdorffraum ist normal \emph{(Übungsaufgabe)}. 
\end{bemerkung}

\begin{definition*}[{name=[Schnitt]}]
Sei $p \colon E \to X$ ein $K$-Vektorbündel über einem topologischen Raum $X$. Ein \Index{Schnitt} von $E$ ist eine stetige Abbildung $s \colon X \to E$ mit 
$p \circ s = \id_X$, das heißt $s(x) \in E_x$. Jedes Vektorbündel hat einen Schnitt, nämlich den trivialen Schnitt $s(x)=0 \in E_x$.	
\end{definition*}

\begin{lemma}[{name=[Fortsetzungen von Schnitten auf abgeschlossenen Teilmengen]},label=schnitte_Hd_Raum]
Sei $X$ ein kompakter Hausdorffraum, $A \subset X$ abgeschlossen. Seien $p \colon E \to X$ und $p' \colon E' \to X$ zwei $K$-Vektorbündel über $X$. Dann gilt:
\begin{enumerate}[a)]
	\item Jeder Schnitt $s \colon A \to E|_A$ lässt sich zu einem Schnitt $\overline{s} \colon X \to E$ fortsetzen.
	\item Sei $\iota \colon A \hookrightarrow X$ die Inklusion und $E|_A := \iota^* E$, sowie $E'|_A := \iota^* E'$. Sei weiter $f \colon E|_A \to E'|_A$ ein Morphismus von 
	Vektorbündeln. Dann existiert ein Morphismus $\hat{f} \colon E \to E'$, der $f$ fortsetzt. Ist $f$ ein Isomorphismus, dann existiert eine offene Menge $U \supset A$, 
	sodass $\hat{f}|_U \colon E|_U \to E'|_U$ ein Isomorphismus von Vektorbündeln ist.
\end{enumerate}	
\end{lemma}

\begin{beweis}
\leavevmode
\begin{enumerate}[a)]
	\item Wir wählen eine endliche Überdeckung $(U_i)_{i=1,\ldots ,N}$ von $X$, sodass lokale Trivialisierungen $\varphi_i\colon E|_{U_i}\to U_i\times K^n$ existieren. Nach 
	dem Fortsetzungssatz \ref{fort_tietze} lässt sich 
	\[
		\pr_{K^n} \circ \varphi_i \circ s\big|_{U_i \cap A} \colon U_i \cap A \longrightarrow K^n 
	\]
	zu einer stetigen Abbildung auf $U_i$ fortsetzen, also auch $s|_{U_i \cap A}$ selbst. Sei $\overline{s}_i$ eine solche Fortsetzung. Sei $\psi_i \colon X \to [0,1]$ eine 
	Partition der Eins\footnote{siehe \url{https://de.wikipedia.org/wiki/Zerlegung_der_Eins}}\index{Partition der Eins}, die der Überdeckung $(U_i)_{i=1,\ldots ,N}$ 
	untergeordnet ist. Sei
	\[
		\overline{s} = \sum_{i=1}^{N} \psi_i \cdot \overline{s}_i \colon X \longrightarrow E 
	\]
	Dann gilt $\overline{s}|_A = \enbrace*{\sum_{i=1}^{N} \psi_i} \cdot s|_A= s|_A$.
	\item Die Morphismen $f \colon E \to E'$ entsprechen den Schnitten des Vektorbündels $\Hom(E,E') \to X$ wie folgt: Jedem $f$ lässt sich ein Schnitt
	\[
		s_f(x)(v) := f(v)
	\]
	für $x \in X$ und $v \in E_x$ zuordnen. Jedem Schnitt $s \colon X \to \Hom(E,E')$ entspricht ein Morphismus $f_s(v)=s \enbrace[\big]{p(v)}(v) \in E'_{p(v)}$. Somit 
	folgt der erste Teil der Aussage aus a).
	
	Sei nun $\pi \colon \Hom(E,E') \to X$ die Bündelabbildung. Sei weiter
	\[
		\Iso(E,E') = \mengenB[\big]{g \in \Hom(E,E')}{g \colon E_{\pi(g)} \to E'_{\pi(g)} \text{ ist ein Isomorphismus}} \subset \Hom(E,E')
	\]
	Dann ist $\Iso(E,E') \to X$ ein lokal triviales Bündel, allerdings \emph{kein} Vektorbündel. Sind $E$ und $E'$ $n$-dimensional, ist die Faser von $\Iso(E,E')$ an jedem 
	Punkt homöomorph zu $\GL_n(K)$. Da $\GL_n(K) \subset M_n(K)$ offen ist und $\Iso(E,E')$ lokal trivial ist, folgt, dass $\Iso(E,E') \subset \Hom(E,E')$ offen 
	ist.\footnote{da lokal homöomorph zu $U \times \GL_n(K) \hookrightarrow U \times M_n(K)$} Sei $\hat{f}$ eine Fortsetzung von $f$. Setze 
	$U = {s_{\hat{f}}}\mkern-0.5mu^{-1} \enbrace[\big]{\Iso(E,E')}$ für $s_{\hat{f}} \colon X \to \Hom(E,E')$. Dann ist $U$ offen und $s_{\hat{f}}|_U$ entspricht einem 
	Isomorphismus $\hat{f} \colon E|_U \to E'|_U$. \qedhere
\end{enumerate}	
\end{beweis}


\begin{lemma}[{name=[{Trivialisierbarkeit von $E \to X \times [a,b]$, wenn $E|_{X \times [a,c]}$, $E|_{X \times [c,b]}$ % 
trivialisierbar}]},label=trivialisierbarkeit_verkleben]
Seien $a,b \in \mathds{R}$ mit $a <b$ und $X$ ein topologischer Raum. Sei weiter $p \colon E \to X \times [a,b]$ ein $K$-Vektorbündel. Es gebe außerdem ein $c \in (a,b)$, 
so dass $E|_{X \times [a,c]}$ und $E|_{X \times [c,b]}$ trivialisierbar sind. Dann ist auch $E$ trivialisierbar.	
\end{lemma}

\begin{beweis}
Wir können den Beweis für jede Komponente von $X$ einzeln führen und daher ohne Beschränkung der Allgemeinheit annehmen, dass $X$ zusammenhängend ist. Seien 
\begin{align*}
	\varphi_a \colon X \times [a,c] \times K^n &\xrightarrow{\minwidthbox{\cong}{2em}} E\big|_{X \times [a,c]} \\ 
	\varphi_b \colon X \times [c,b] \times K^n &\xrightarrow{\minwidthbox{\cong}{2em}} E\big|_{X \times [c,b]}
\end{align*}
entsprechende Trivialisierungen. Sei außerdem
\[
	h := \enbrace*{\varphi_b\big|_{X \times \set{c} \times K^n }} ^{-1} \circ  \enbrace*{\varphi_a\big|_{X \times \set{c} \times K^n}} \colon X \times \set{c} \times K^n 
	\longrightarrow X \times \set{c} \times K^n    
\]
Dann ist $h$ ein Isomorphismus eines trivialen $K$-Vektorbündels über $X$.
Dieser hat die Form $h(x,v)= \enbrace[\big]{x , g(x) \cdot v}$ für eine stetige Abbildung $g \colon X \to \GL_n(K)$. Sei jetzt 
\[
	w \colon X \times [c,b] \times K^n \longrightarrow X \times [c,b] \times K^n \quad , \qquad (x,t,v') \longmapsto \enbrace[\big]{x,t,g(x)\cdot v'}
\]
Dann gilt $\varphi_a|_{X \times \set{c} \times K^n} = \varphi_b \circ w |_{X \times \set{c} \times K^n }$ und $w$ ist ein Isomorphismus. Daher können wir diese beiden 
Abbildungen zu einem Isomorphismus von Vektorbündeln $X \times [a,b] \times K^n \to E$ zusammensetzen.
\end{beweis}

\begin{lemma}[{name=[{Endliche Überdeckung $\set{U_i}_{i=1}^N$, sodass Trivialisierbarkeit auf $U_i$ gilt}]},label=trivialisierbarkeit_ueberdeckung]
Sei $X$ ein kompakter Raum. Sei $p \colon E \to X \times I$ ein $K$-Vektorbündel. Dann gibt es eine endliche offene Überdeckung $\set{U_i}_{i=1}^N$ von $X$, so dass 
$E|_{U_i \times I}$ trivialisierbar ist.	
\end{lemma}

\begin{beweis}
Sei $x \in X$ und $t \in I$. Dann existieren offene Umgebungen $x \in U(t) \subset X$ und $I(t) \subset I$, sodass $E|_{U(t) \times I(t)}$ trivialisierbar ist. Da $I$ 
kompakt ist, existiert eine endliche Folge $0=t_0 < t_1 < \ldots < t_n =1$ zusammen mit offenen Umgebungen $U_i$ von $x$, so dass $E|_{U_i \times [t_{i-1}, t_i]}$ 
trivialisierbar ist. Sei nun $U := \bigcap_{i=1}^n U_i$. Nach \autoref{trivialisierbarkeit_verkleben} ist $E|_{U \times I}$ trivialisierbar. Da $X \times I$ kompakt ist, 
wird es von endlichen vielen solcher Mengen überdeckt. 	
\end{beweis}


\begin{theorem}[{name=[{Isomorphismus von Vektorbündeln $E$ und $E|_{X \times \set{1}} \times I$ für $X$ kompakt}]},label=iso_vektb_intervall]
Sei $X$ ein kompakter Raum und sei $p \colon E\to X \times I$ ein $K$-Vektorbündel. Sei weiter $E_1 = E|_{X \times \set{1}}$.
\begin{enumerate}[a)]
	\item Es gibt einen Isomorphismus von Vektorbündeln über $X \times I$
	\[
		\begin{tikzcd}
			\varphi \colon E \rar["\cong"]  & E_1 \times I 
		\end{tikzcd}
	\]
	mit $\varphi|_{p ^{-1}(X \times \set{1})} =\id_{E_1}$
	\item Sei $(X,A)$ ein \hyperref[def_koraumpaar]{Ko-Raumpaar}, $h \colon X \times I \to Y$ eine stetige Abbildung und $p' \colon E' \to Y$ ein $K$-Vektorbündel mit $E \cong h^* E'$. Falls jetzt
	$h_t(a)=h_0(a)$ für alle $a \in A$ und $t \in I$, dann lässt sich der Isomorphismus aus a) so wählen, dass er den durch $h$ induzierten Isomorphismus 
	$E|_{A \times I} \to E_1|_{A} \times I$ fortsetzt.
\end{enumerate}	
\end{theorem}

\begin{beweis}
Nach Lemma \ref{trivialisierbarkeit_ueberdeckung} können wir eine Überdeckung von $X \times I$ durch offene Mengen der Form $U_i \times I$ mit $i \in \set{1,\ldots ,N}$ 
finden, so dass $E|_{U_i \times I}$ trivialisierbar ist. Seien $h_i \colon U_i \times I \times K^n \to E|_{U_i \times I}$ die entsprechenden Trivialisierungen und seien 
$\eta_i \colon X \to [0,1]$ stetige Funktionen mit folgenden Eigenschaften:
\[
	\supp(\eta_i) \subset U_i \qquad \text{ und } \qquad \max_{i=1,\ldots ,N} \eta_i(x)=1 \text{ für alle } x \in X
\]
Sei $r_i \colon X \times I \to X \times I$ gegeben durch $r_i(x,t)= \enbrace[\big]{x, \max(\eta_i(x),t)}$. Ein Punkt $(x,t)$ wird durch $r_i$ entlang der $I$-Achse nach $1$ verschoben. Der 
Wert der Verschiebung wird durch $\eta_i$ festgelegt.
\[
	\begin{tikzpicture}[my brace/.style={decorate, decoration={brace, amplitude=5pt, mirror, raise=5pt}}]
		\draw[thick] (0,0) node[below]{$0$} rectangle (2,2);
		\draw[|-|] (0,1.6) -- (0,0.5);
		\draw[my brace,] (0,1.6) -- (0,0.5) node[midway, left=8pt]{$U_i$};
		\foreach \x in {0.3,0.65,...,1.8}{
			\draw[gray] (\x,0.2) -- (\x,1.8);
		}
		\node[below] at (2,0) {$1$};
		
		\draw[thick,-to] (2.5,1) -- (4,1) node[midway,above]{$r_i$};
		
		
		\node[below] at (7,0) {$1$};
		\draw[|-|] (5,1.6) -- (5,0.5);
		\draw[my brace,] (5,1.6) -- (5,0.5) node[midway, left=8pt]{$U_i$};
		\foreach \x in {5.3,5.65,...,6.8}{
			\draw[gray] (\x,0.2) -- (\x,0.5) .. controls +(90:0.2) and +(270:0.2) .. (7,1.05) .. controls +(90:0.2) and +(270:0.2) .. (\x,1.6) -- (\x,1.8);
		}
		\draw[gray,->] (5.5,1.05) -- ++(0.7,0);
		\draw[thick] (5.,0) node[below]{$0$} rectangle (7,2);
	\end{tikzpicture}
\]
Nun konstruieren wir einen Morphismus von $K$-Vektorbündeln $(u_i,r_i)$ über $r_i$ wie folgt: 
\[
	u_i \enbrace[\big]{h_i(x,t,v)} := h_i \enbrace[\big]{r_i(x,t),v}  
\]
für $(x,t,v) \in U_i \times I \times K^n$. Wir setzen $u_i=\id$ außerhalb des Bildes von $h_i$. Wegen 
$r_i|_{X \times I\setminus{U_i} \times I} =\id_{X \times I \setminus U_i \times I}$ ist dies stetig. Außerdem kommutiert
\[
	\begin{tikzcd}
		E \rar["u_i"] \dar["p"]& E \dar["p"]\\
		X\times I \rar["r_i"] & X \times I
	\end{tikzcd}
\]
und $u_i$ ist auf den Fasern linear. Sei $r :=r_N \circ r_{N-1} \circ \ldots \circ r_1$. Nach Konstruktion gilt $r(x,t)=(x,1)$ für alle $x \in X$. 
Sei $u := u_N \circ \ldots \circ u_2 \circ u_1$.Dann ist $(u,r)$ ein Morphismus von Vektorbündeln und es gilt $(u,r)|_{X \times \set{1}}=\id_{E_1}$. Außerdem haben wir 
\[
	\begin{tikzcd}
		E \rar["\enbrace*{u,r} "] \drar & E \\
		& E_1 \uar[hook]
	\end{tikzcd}
	\qquad \text{ und }\qquad 
	\begin{tikzcd}
		E \rar["u"] \dar["p"'] & E_1 \dar["p_1"] \\
		X \times I \rar["r"] & X \times \set{1} 
	\end{tikzcd}
\]
Setze jetzt $\varphi \colon E \to E_1 \times I$, $v \mapsto \enbrace[\big]{u(v), \pr_I \circ p(v)} $. Dann ist $(\varphi, \id_{X \times I})$ ein Isomorphismus von 
Vektorbündeln mit $\varphi|_{E \times \set{1}} =\id$. Damit ist a) gezeigt.

Für b) sei $V_j \subset Y$ eine offene Überdeckung von $Y$, sodass $E'|_{\overline{V}_j}$ trivialisierbar ist. Sei nun $B_j = h ^{-1}(\overline{V}_j) \subset X \times I$. 
Es gilt dann
\[
	B_j \cap (A \times I) = \enbrace[\big]{\Underbrace{h_1 ^{-1}(\overline{V}_j) \cap A}{=A_j}} \times I
\]
da $h$ auf $A$ unabhängig ist von $t \in I$. Nach Konstruktion ist $E|_{A_j \times I}$ trivialisierbar. Nach \autoref{schnitte_Hd_Raum} b) finden wir eine offene Umgebung 
$U_j \times I$ mit $X \supset U_j \supset A_j$, so dass $E|_{U_j \times I}$ trivialisierbar ist. Wir ergänzen die Mengen $U_j \times I$ um trivialisierende offene Umgebungen 
$U_i \times I$ mit $U_i \cap A = \emptyset$, so dass $\bigcup_{i} U_i \times I = X \times I$. Jetzt verwenden wir die Konstruktion aus a) für diese Überdeckung. Der 
entsprechende Isomorphismus $\varphi \colon E \to E_1 \times I$ stimmt auf $A_j \times I$ mit dem von $h$ induzierten Isomorphismus überein. 	
\end{beweis}
% subsection 210 (end)

Das folgende Korollar ist von entscheidender Bedeutung für die Definition von topologischer $K$-Theorie, die wir im nächsten Abschnitt erarbeiten wollen.

\begin{korollarB}[{name=[Homotopieinvarianz von Pullbacks]}]
Sei $X$ ein kompakter Raum und $h \colon X \times I \to Y$ eine Homotopie. Sei $p \colon E \to Y$ ein $K$-Vektorbündel. Dann gilt
\begin{enumerate}[a)]
	\item $h_0^* E \cong h_1^* E$ als Vektorbündel über $X$, das heißt Pullbacks sind homotopieinvariant bis auf Isomorphie.
	\item Falls $(X,A)$ ein Ko-Raumpaar ist und $h$ stationär auf $A$ ist, dann kann der Isomorphismus so gewählt werden, dass er über $A$ die Identität ist.\qedhere
\end{enumerate}	
\end{korollarB}

% section 2 (end)
\newpage

\section{Topologische $K$-Theorie} % (fold)
\label{sec:3}
Wir haben im letzten Kapitel folgenden kontravarianten, homotopieinvarianten Funktor kennengelernt
\[
	\begin{tikzcd}[row sep=0.1em, column sep=large,cramped]
		\TOP \rar & \MONOIDE \\
		X \rar[mapsto] & \nicefrac{\VEKT_K(X)}{\sim_{\mathrm{Iso}}} \\
		(f \colon X \to Y) \rar[mapsto] & \enbrace[\big]{f^* \colon \nicefrac{\VEKT_K(Y)}{\sim_{\mathrm{Iso}}} \to \nicefrac{\VEKT_K(X)}{\sim_{\mathrm{Iso}}}} 
	\end{tikzcd}
\]
mit der Monoid-Operation $[E]_{\Iso} + [F]_{\Iso} = [E \oplus F]_{\Iso}$ und dem Nullvektorbündel als neutralem Element. Daraus wollen wir eine Gruppe konstruieren. Ein kanonischer Ansatz dazu ist die folgende Konstruktion: 

\subsection*{Grothendieck-Konstruktion}
Sei $(M,+)$ ein kommutativer Monoid. Betrachte $\Gr(M)= \nicefrac{M \times M}{\sim}$ mit der Äquivalenzrelation
\[
	(x,y) \sim (x',y') :\iff \exists c \in M \text{ mit } x+y' +c = x' +y +c
\]
Die Idee dahinter: $[x,y]$\enquote{$=$}$x-y$, also die erste Koordinate ist der \enquote{positive} und die zweite der \enquote{negative} Teil. Dann ist 
\[
	\begin{tikzcd}[row sep=tiny]
		+ \colon \Gr(M) \times \Gr(M) \rar & \Gr(M) \\
		\quad \enbrace[\big]{[x,y] , [x',y']} \rar[mapsto] & \benbrace*{x+x',y+y'}  
	\end{tikzcd}
\]
wohldefiniert. Mit dieser Operation ist $\Gr(M)$ eine abelsche Gruppe (das Inverse zu $[x,y]$ ist $[y,x]$, denn $[x+y,x+y]=[0,0]$) und $M \to \Gr(M)$, $x  \mapsto [x,0]$ ist
ein Monoidhomomorphismus.
\begin{beispiel}
$\Gr(\mathds{N}_0,+)= (\mathds{Z},+)$.
\end{beispiel}
\minisec{Universelle Eigenschaft}
Zu jedem Monoidhomomorphismus $\varphi \colon M \to A$ in eine abelsche Gruppe $A$ existiert genau ein Gruppenhomomorphismus $\psi \colon \Gr(M) \to A$, sodass das Diagramm
\[
	\begin{tikzcd}
		M \dar[hook] \rar["\varphi"] & A \\
		\Gr(M) \urar["\exists ! \psi"', dashed]
	\end{tikzcd}
\]
kommutiert. \emph{Der Beweis ist eine Übungsaufgabe. Für die Existenz betrachte $\psi \enbrace*{[x,y]} = \varphi(x) - \varphi(y)$. Die Eindeutigkeit folgt dann ganz 
einfach.} \bigskip \\
Wir bekommen jetzt eine abelsche Gruppe $K^0(X) := \Gr \enbrace*{\nicefrac{\VEKT_{\mathds{C}}(X)}{\sim_{\Iso}}}$, aber es stellen sich weiterhin die Fragen auf dem Weg zu 
einer Kohomologie-Theorie: Was ist $K^n(X)$ für $n\not=0$? Was ist $K^n(X,A)$, also was machen wir mit (Ko-)Raumpaaren?

\begin{definition}[{name=[Tripel, Morphismus von Tripeln und stabile Isomorphie]}, label=sub:31]
Sei $(X,A)$ ein \hyperref[def_koraumpaar]{Ko-Raumpaar}.
\begin{enumerate}[(i)]
	\item Ein \Index{Tripel} $(E_1,f,E_0)$ über $(X,A)$ besteht aus endlichdimensionalen $\mathds{C}$-Vektorbündeln $p_i \colon E_i \to X$ und einem Isomorphismus
	$f \colon E_1|_A \to E_0|_A$.
	\item Ein \Index{Morphismus von Tripeln} $(E_1,f,E_0)$ und $(E_1',f',E_0')$ besteht aus einem Paar $(g_1,g_0)$ von Vektorbündelmorphismen $g_i \colon E_i \to E_i'$,
	sodass das folgende Diagramm kommutiert
	\[
		\begin{tikzcd}[column sep=large]
			E_1|_A \rar["g_1|_A"] \dar["f"]& E_1'|_A \dar["f'"]\\
			E_0|_A \rar["g_0|_A"] & E_0'|_A
		\end{tikzcd}
	\] 
	\item Tripel der Form $(E,\id_E,E)$ heißen \bet{elementar}\index{Tripel!elementares}.
	\item $(E_1,f,E_0)$ und $(E_1',f',E_0')$ heißen \Index{stabil isomorph}, falls Vektorbündel $E \to X$ und $E' \to X$ existieren, so dass
	\[
		\enbrace*{E_1 \oplus E, f \oplus \id_E, E_0 \oplus E} \cong \enbrace*{E_1' \oplus E', f' \oplus \id_{E'}, E_0' \oplus E' } 
	\]
	also falls sie isomorph werden nach Addition elementarer Tripel.
\end{enumerate}
\end{definition}
\begin{bemerkung}
\begin{itemize}
	\item Stabile Isomorphie ist eine Äquivalenzrelation.\hfill \emph{(einfache Übung)}
	\item Stabile Isomorphieklassen von Tripeln lassen sich addieren, das heißt
	\[
		\benbrace[\big]{E_1,f,E_0}_\st + \benbrace[\big]{E_1',f',E_0'}_\st := \benbrace[\big]{E_1 \oplus E_1', f \oplus f', E_0 \oplus E_0'}_\st  
	\]
	ist wohldefiniert. Das Tripel $\benbrace*{0,\id,0}_\st$ ist ein neutrales Element hierfür.
	\item Es gilt $\benbrace*{E_1,f,E_0}_\st=\benbrace*{0,\id,0}_\st$ genau dann, wenn ein Vektorbündel $E \to X$ existiert, so dass sich 
	der Isomorphismus $f \oplus \id \colon E_1  \oplus E|_A \to E_0 \oplus E|_A $ zu einem Isomorphismus über ganz $X$ ausdehnen lässt.
\end{itemize}
\end{bemerkung}
\begin{beweis}
Für die erste Implikation folgt aus der Definition $(E_1 \oplus E, f \oplus \id, e_0 \oplus E) \cong \enbrace*{E', \id_{E'},E'}$. Sei $\varphi=(\varphi_1,\varphi_0)$ dieser 
Isomorphismus. Dann kommutiert das folgende Diagramm
\[
	\begin{tikzcd}[column sep=large]
		(E_1 \oplus E)|_A \rar["\varphi_1|_A","\cong"'] \dar["f \oplus \id"']& E'|_A \dar["\id_{E^\prime}\!|_A","\cong"'] \\
		(E_0 \oplus E)|_A \rar["\cong","\varphi_0|_A"'] & E'|_A
	\end{tikzcd}
\]
Der Isomorphismus $\varphi_0 ^{-1} \circ \varphi_1 \colon E_1 \oplus E \to E_0 \oplus E$ ist auf ganz $X$ definiert und erfüllt 
$\varphi_0 ^{-1} \circ \varphi_1 |_A = f \oplus \id$, ist also eine Fortsetzung wie gewünscht.
% $\id_{E'}$ ist ein Isomorphismus auf ganz $X$, also lässt sich $f \oplus \id$ zu einem Isomorphismus auf ganz $X$ ausdehnen.

Für die Rückrichtung betrachten wir $\psi \colon E_1 \oplus E \xrightarrow{\minwidthbox{\raisebox{0pt}[0pt][0pt]{$\cong$}}{2em}} E_0 \oplus E$ mit $\psi|_A = f \oplus \id$. Dann kommutiert
\[
	\begin{tikzcd}[column sep=large]
		(E_1 \oplus E)|_A  \dar["f \oplus \id"] \rar["\psi|_A","\cong"']& (E_0 \oplus E)|_A \dar["\id"] \\
		(E_0 \oplus E)|_A \rar["\id","\cong"']& (E_0 \oplus E)|_A
	\end{tikzcd} 
\]	
also ist $(\psi,\id) \colon \enbrace*{E_1 \oplus E, f \oplus \id, E_0 \oplus E} \to \enbrace*{E_0 \oplus E, \id ,E_0 \oplus E}$ die gesuchte stabile Isomorphie.
\end{beweis}
% subsection 31 (end)


\begin{definition}[{name=[$K^{-n}(X,A)$]},label=sub:32]
Sei $(X,A)$ ein \hyperref[def_koraumpaar]{Ko-Raumpaar}. Wir definieren $K^0(X,A)$ als die Menge der stabilen Isomorphieklassen von Tripeln über $(X,A)$. 
$K^0(X,A)$ heißt auch \emph{$0$-te $K$-Gruppe} von $(X,A)$. Für $n>0$ definieren wir \index{K-Gruppe@$K$-Gruppe}
\begin{align*}
	K^{-n}(X,A) &:= K^0 \enbrace[\big]{(X,A) \times (I,\partial I)^n} \\
	\text{und } \quad K^{-n}(X)\hphantom{,A} &:= K^{-n}(X,\emptyset)
\end{align*}
\end{definition}
\noindent \emph{Übung: $K^0(X,\emptyset)$ ist isomorph zu $\Gr \enbrace*{\nicefrac{\VEKT_\mathds{C}(X)}{\sim_{\Iso}}}$.} 

Wir wissen bisher nur, dass $K^0(X,A)$ ein Monoid ist bezüglich der Addition von Tripeln. Um zu sehen, dass dies tatsächlich eine Gruppenstruktur liefert, brauchen wir zwei 
Lemmata: 
% subsection 32 (end)


\begin{lemma}[label=sub:33,{name=[Existenz eines Bündelisomorphismus und einer Homotopie]}]
Sei $(X,A)$ ein Ko-Raumpaar.\marginnote{Morphismen entsprechen Schnitten von Morphismenbündeln, siehe Beweis von \autoref{schnitte_Hd_Raum}}
Seien $p \colon E \to X$ und $p' \colon E' \to X$ zwei $K$-Vektorbündel und $f \colon E'|_A \to E|_A$ ein Isomorphismus, $(E',f,E)$ also ein Tripel. 
Sei weiter $f' \colon E' \to E$ ein Isomorphismus von $K$-Vektorbündeln und $h \colon A \times I \to \Iso(E'\times I|_{A\times I},E\times I|_{A\times I})$ eine stetige 
Abbildung, so dass $h_t$ ein Schnitt ist für alle $t \in I$ und $h_0=f'|_A$, sowie $h_1=f$ gilt.\todo[inline]{Ich sehe nicht ein, wie bei dieser Definition von $h$, $h_1=f$ bzw. $h_0=f'|_A$ gelten kann, da in deren Definitionen das Einheitsintervall doch gar nicht vorkommt \ldots}

Dann existiert ein Isomorphismus $F \colon E' \to E$ mit $F|_A=f$ und eine Homotopie $H \colon X \times I \to \Iso(E' \times I,E\times I)$ mit $H_0=f'$, $H_1=F$ und $H|_{A\times I}=h$.	
\end{lemma}
\begin{beweis}
Unser Ziel ist es $h$ zu einer Homotopie $H$ auf ganz $X \times I$ zu erweitern und dann $F = H_1$ zu setzen.
Die Abbildungen $h$ und $f'$ liefern eine stetige Abbildung 
\[
	\begin{tikzcd}
		\tilde{h} \colon A \times I \cup X \times \set{0} \rar & \Iso \enbrace*{\enbrace*{\pr_X^* E'}|_{A \times I \cup X \times \set{0} },\enbrace*{\pr_X^* E}|_{A \times I \times X \times 
		\set{0}}} 
	\end{tikzcd}
\]
Wir müssen $\tilde{h}$ zu einem Schnitt $X \times I \to \Iso(E'\times I, E \times I)$ ausdehnen, dass wir einen Schnitt bekommen. Betrachte \marginnote{$r$ aus Kofaserungseigenschaft}
\[
	\begin{tikzcd}
		A \times I \cup X \times \set{0} \rar["\id"] \dar[hook]& A \times I \cup X \times \set{0} \rar["\iota",hook] & X \times I\\  
		X \times I \urar[dashed,"r"]
	\end{tikzcd}
\]
Sei $r_X = \pr_X \circ \iota \circ r$ und $r_I = \pr_I \circ \iota \circ r$. Sei weiter 
\[
	\begin{tikzcd}[row sep=tiny]
		K \colon (X \times I) \times I \rar &  X \times I \\
		(x,t,s) \rar[mapsto] & \enbrace*{r_X(x,st), s r_I(x,t) + (1-s)t} 
	\end{tikzcd}
\]
Dann gilt $K_0(x,t)=(x,t)$ und $K_1(x,t)= \iota \circ r(x,t)$, also ist $\iota \circ r$ homotop zu $\id_{X \times I}$. Folglich ist
\[
	\begin{tikzcd}
		E \times I \rar["\cong"] &(\iota \circ r)^*(E \times I) \cong r^* \enbrace*{\iota^*(E \times I)} \rar["\overline{r}"] & \iota^*(E \times I) \cong E|_A \times I \cup E 
		\times \set{0} 
	\end{tikzcd}
\]
Ebenso $\overline{r}' \colon E' \times I E'|_A \times I \cup E' \times \set{0}$. $(\overline{r},r)$ und $(\overline{r}',r)$ sind Morphismen von Vektorbündeln.
Wir erhalten nun das folgende Diagramm, wobei die diagonale Abbildung rechts oben von $\tilde{h}$ induziert ist und somit ein Isomorphismus ist
\[
	\begin{tikzcd}[] % row sep=2.3em, column sep=2.8em
		E' \times I \ar[dd] \ar[rr,"\overline{r}'"] \ar[dr,"\overline{H}",dashed]& & E'|_A \times I \cup E' \times \set{0} \ar[dd] \drar["\cong"'] \\
		& E \times I \ar[rr,"\overline{r}", crossing over,pos=0.3] & & E|_A \times I \cup E \times \set{0} \ar[dd] \\
		X \times I \ar[rr,"r",pos=0.3] \ar[dr,"\id"] & & A \times I \cup X \times \set{0} \ar[dr,"\id"] \\
		& X \times I \ar[from=uu,crossing over] \ar[rr,"r"']& & A \times I \cup X \times \set{0} 
	\end{tikzcd}
\]
Die Existenz von $\overline{H}$ folgt aus der universellen Eigenschaft des Pullbacks. Es gilt nun $H=\pr_E \circ \overline{H} \colon E' \times I \to E$ und wir erhalten
eine Abbildung $X \times I \to \Iso(E' \times I, E \times I)$. $F := H|_{E' \times \set{1}} \colon E' \to E$ ist nun die gesuchte Abbildung.
\end{beweis}
% subsection 3.3 (end)

Wir können nur klären, was inverse Elemente in $K^0(X,A)$ sind. Kandidat für das inverse Element in $K^0(X,A)$ zu $(E_0,f,E_1)$ ist $(E_1, f ^{-1},E_0)$. Wir müssen
prüfen, dass $(E_0,f,E_1) \oplus (E_1,f ^{-1}, E_0)= (E_0 \oplus E_1, f \oplus  f ^{-1},E_1 \oplus E_0)$ ist. Dazu folgendes Lemma:

\begin{lemma}[label=sub:34]
Sei $(X,A)$ ein Ko-Raumpaar.
\begin{enumerate}[(i)]
	\item Seien $(E_0,f,E_1)$ und $(E_1,g,E_2)$ Tripel über $(X,A)$. Dann sind die beiden Tripel 
	\[
		(E_0,g \circ f, E_2) \quad  \text{ und } \quad \enbrace*{E_0 \oplus E_1, f \oplus g, E_1 \oplus E_2} 
	\]
	stabil isomorph.
	\item Seien $(E_0,f,E_1)$ und $(E_0,g,E_1)$ Tripel über $(X,A)$. Sei $h \colon A \times I \to \Iso \enbrace*{E_0|_A,E_1|_A} $ eine Homotopie zwischen $f$ und $g$, sodass 
	$h_t \colon A \to \Iso \enbrace*{E_0|_A,E_1|_A}$ ein Schnitt ist für alle $t \in I$. Dann sind die beiden Tripel stabil isomorph.
\end{enumerate}
\end{lemma}
\begin{beweis}
\mbox{ }
\begin{enumerate}[(i)]
	\item Wir haben das folgende kommutative Diagramm von Vektorbündeln:
	\[
		\begin{tikzcd}[ampersand replacement=\&,row sep=4.5em,column sep=5.2em]
			E_1 \oplus E_0\big|_A \ar[rrr,"{\begin{pmatrix} \id & f \\ 0 & \id \end{pmatrix}}","\text{(1)}"'] \ar[d,"{\begin{pmatrix} \id & 0 \\ 0 & g \circ f \end{pmatrix}}"'] 
			\& \& \& E_1 \oplus E_0\big|_A \ar[d,"{\begin{pmatrix} g & 0 \\ 0 & f \end{pmatrix}}"] \\
			E_1 \oplus E_2\big|_A \& E_2 \oplus E_1\big|_A \lar["{\begin{pmatrix} 0 & -\id \\ \id & 0 \end{pmatrix}}"',"\text{(2)}"]
			\& E_2 \oplus E_1\big|_A \lar["{\begin{pmatrix} \id & g \\ 0 & \id \end{pmatrix}}"',"\text{(3)}"]
			\& E_2 \oplus E_1\big|_A \lar["{\begin{pmatrix} \id & 0 \\ - g^{-1} & \id \end{pmatrix}}"',"\text{(4)}"]
		\end{tikzcd}
	\]
	Die Kommutativität ergibt sich dabei wie folgt
	\begin{align*}
		&\begin{pmatrix} 0 & -\id \\ \id & 0 \end{pmatrix} \cdot 
		\begin{pmatrix} \id & g \\ 0 & \id \end{pmatrix} \cdot 
		\begin{pmatrix} \id & 0 \\ - g^{-1} & \id \end{pmatrix} \cdot 
		\begin{pmatrix} g & 0 \\ 0 & f \end{pmatrix} \cdot 
		\begin{pmatrix} \id & f \\ 0 & \id \end{pmatrix} \\
		&= \begin{pmatrix} 0 & -\id \\ \id & 0 \end{pmatrix} \cdot 
		\begin{pmatrix}	0 & g \\ - g^{-1} & \id \end{pmatrix} \cdot 
		\begin{pmatrix} g & g \circ f \\ 0 & f \end{pmatrix} = 
		\begin{pmatrix}	g ^{-1} & -\id \\ 0 & g \end{pmatrix} \cdot 
		\begin{pmatrix} g & g \circ f \\ 0 & f \end{pmatrix} = 
		\begin{pmatrix}\id & 0 \\ 0 & g \circ f \end{pmatrix}
	\end{align*}
	Wir können $f$,$g$ und $g^{-1}$ zu Morphismen über $X$ fortsetzen (nach \autoref{schnitte_Hd_Raum}). Die Fortsetzungen sind im Allgemeinen aber keine Isomorphismen.
	Dennoch sind die auf $X$ erweiterten Versionen der Morphismen (1) -- (4) Isomorphismen über $X$. Dabei ist zum Beispiel das Inverse zu $\begin{psmallmatrix}
		\id & f \\ 0 &\id
	\end{psmallmatrix}$ gegeben durch $\begin{psmallmatrix}
		\id & -f \\ 0 & \id
	\end{psmallmatrix}$. Dass die beiden Tripel stabil isomorph sind, ergibt sich nun wie folgt:
	\begin{align*}
		\benbrace[\big]{E_0 , g \circ f, E_2}_\st =
		\benbrace[\big]{E_1 \oplus  E_0, \id \oplus g \circ f, E_1 \oplus E_2}_{\st} &= \benbrace*{E_1 \oplus E_0, \begin{psmallmatrix}
			g & 0 \\ 0 &f
		\end{psmallmatrix}, E_2 \oplus E_1}_{\st}  \\ 
		&= \benbrace*{E_0 \oplus E_1, \begin{psmallmatrix}
			f & 0 \\ 0 &g
		\end{psmallmatrix}, E_1 \oplus E_2}_{\st} 
	\end{align*}
	\todo{Revision}
	\item Betrachte das Tripel $(E_0, \oplus E_1, f \oplus g ^{-1}, E_1 \oplus E_0)$. Nach (i) ist es stabil isomorph zu $(E_0, g ^{-1}\circ f,E_0)$. Nach Voraussetzung ist
	$g ^{-1} \circ f$ homotop zu $\id_{E_0}$ durch Isomorphismen. Also folgt aus Lemma \ref{sub:33}, dass sich $g ^{-1} \circ f$ zu einem Isomorphismus über $X$ fortsetzen 
	lässt. Aus der Bemerkung folgt 
	\begin{align*}
		\benbrace*{E_0, g ^{-1} \circ f,E_0}_\st = \benbrace*{0,\id,0}_\st 
	\end{align*}
	Außerdem gilt nach (i) auch $\benbrace*{E_1 \oplus E_0, g ^{-1} \oplus g, E_0 \oplus E_1}_\st = \benbrace*{0,\id,0}$. Weiter ist
	\begin{align*}
		\benbrace*{E_0,f,E_1}_\st &= \benbrace*{E_0 \oplus E_1 \oplus E_0, f \oplus g ^{-1} \oplus g, E_1 \oplus E_0 \oplus E_1}_\st \\  
		&= \Underbrace{\benbrace*{E_0 \oplus E_1, f \oplus g ^{-1}, E_1 \oplus E_0}_\st }{=\benbrace{0,\id,0}} + \benbrace*{E_0,g,E_1}_\st \qedhere
	\end{align*}
\end{enumerate}
\end{beweis}
% subsection 34 (end)

\subsection{Korollar} % (fold)
\label{sub:35}
\minisec{Beweis}
Es genügt zu zeigen, dass $K^0(X,A)$ eine abelsche Gruppe ist. Sei $\benbrace*{E_0, f,E_1}_\st\in K^0(X,A)$. Dann gilt nach Lemma \ref{sub:34}
\begin{align*}
	\benbrace*{E_0,f,E_1}_\st + \benbrace*{E_1,f ^{-1},E_0}_\st = \benbrace*{E_0 \oplus E_1, f \oplus f ^{-1}, E_1 \oplus E_0}_\st = \benbrace*{E_0, f ^{-1} \circ f,E_0}_\st
	= \benbrace*{0,\id,0}_\st \bewende  
\end{align*}
Die Zuordnung $(X,A) \mapsto K^{-n}(X,A)$ ist ein kontravarianter Funktor: Sei $g \colon (X,A) \to (Y,B)$ eine stetige Abbildung von Ko-Raumpaaren, dann ist
\[
	g^* \colon K^{-n}(Y,B) \to K^{-n}(X,A)
\]
definiert durch $g^* \benbrace*{E_0,f,E_1}_\st = \benbrace*{g^* E_0, g^* f, g^* E_1}_\st$. Hierbei ist $g^* f$ der von der universellen Eigenschaft des Pullbacks induzierte 
Morphismus:
\[
	\begin{tikzcd}[]
		g^* E_0\big|_A \ar[rr] \ar[dd] \ar[dr,dashed,"g^* f","\cong"'] & & E_0\big|_B \drar["f","\cong"'] \ar[dd]\\
		& g^* E_1\big|_A \ar[rr,crossing over] & & E_1\big|_B \ar[dd]\\
		A \ar[rr,"g|_A",pos=0.3] \ar[dr,"\id"']& & B \ar[dr,"\id"]\\
		& A \ar[from=uu,crossing over] \ar[rr,"g|_A"]& & B
	\end{tikzcd}
\]
% subsection 35 (end)

\subsection{Lemma} % (fold)
\label{sub:36}
Der Funktor $(X,A) \mapsto K^{-n}(X,A)$ ist homotopieinvariant.
\minisec{Beweis}
Es genügt wieder, den Fall $n=0$ zu betrachten. Seien $g_0,g_1 \colon (X,A) \to (Y,B)$ und sei $h \colon (X \times,A \times I) \to (Y,B)$ eine Homotopie zwischen $g_0$ und 
$g_1$. Sei $(E_0,F,E_1)$ ein Tripel über $(Y,B)$. Nach Theorem \ref{iso_vektb_intervall} a) existieren Isomorphismen
\[
	\begin{tikzcd}[column sep=large, row sep=0.5em]
		h^* E_0 \rar["\psi^\ssbrace{0}","\cong"'] & g_1^*(E_0) \times I \\
		h^* E_1 \rar["\psi^\ssbrace{1}","\cong"'] & g_1^*(E_1) \times I
	\end{tikzcd}
\]
die über $X \times \set{1}$ die Identität sind. Betrachte 
\[
	\enbrace*{\psi^\ssbrace{1}\big|_{A \times I}} \circ \enbrace*{h^* f} \circ \enbrace*{\psi^\ssbrace{0}\big|_{A \times I}} ^{-1} \colon g^*_1 E_0\big|_A \times I  
	\xrightarrow{\cong} g_1^*E_1\big|_A \times I   
\]
Dies entspricht einer Homotopie $\tilde{h} \colon A \times I \to \Iso \enbrace*{g_1^*E_0\big|_A, g_1^* E_1\big|_A}$. Sei $\tilde{f} = \enbrace*{\psi^\ssbrace{1}\big|_{A \times \set{0}}} \circ \enbrace*{h^*_0 f} \circ \enbrace*{\psi^\ssbrace{0}\big|_{A \times \set{0}}} ^{-1} $. Dann haben wir das kommutative Diagramm
\[
	\begin{tikzcd}[sep=large]
		g_0^* E_0\big|_A \rar["g_0^*f","\cong"'] \dar["\psi^\ssbrace{0}|_{A \times \set{0}}"',"\cong"] & g_0^*E_1\big|_A \dar["\psi^\ssbrace{1}|_{A \times \set{0} }","\cong"'] \\
		g^*_1 E_0\big|_A \rar["\cong"',"\tilde{f}"] & g_1^* E_1\big|_A
	\end{tikzcd}
\]
Es folgt damit
\begin{align*}
	\benbrace[\big]{g_0^* E_0, g^*_0 f , g^*_0 E_1}_\st = \benbrace*{g_1^* E_0, \tilde{f}, g_1^* E_1}_\st \stackrel{\mbox{\scriptsize \ref{sub:34} ii)}}{=}
	\benbrace[\big]{g_1^* E_0, g_1^* f, g_1^* E_1}_\st \bewende 
\end{align*}
% subsection 36 (end)

\subsection*{Ausschneidung}
Sei $(X,A)$ ein Ko-Raumpaar und $U \subset A$, sodass $\overline{U} \subset \mathring{A}$, dann ist die induzierte Abbildung 
$K^{-n}(X,A) \to K^{-n}(X \setminus U, A \setminus U)$ ein Isomorphismus.

\subsection{Lemma} % (fold)
\label{sub:37}
Sei folgendes Diagramm ein Pushout-Diagramm von kompakten Hausdorffräumen
\[
	\begin{tikzcd}
		X_0 \rar[hook,"i_1"] \dar["i_2"]& X_1 \dar["j_1"] \\
		X_2 \rar[hook,"j_2"] & X
	\end{tikzcd}
\]
Seien $p_i \colon E_i \to X_i$ für $i \in \set{1,2}$ $K$-Vektorbündel und $f \colon i_1^* E_1 \xrightarrow{\cong} i_2^* E_2 $ ein Isomorphismus von $K$-Vektorbündeln über 
$X_0$. Dann ist \todo{Hier vlt mal \texttt{xfrac} benutzen?}
\[
	E = \nicefrac{E_1 \coprod E_2}{\sim} \xrightarrow{\minwidthbox{p}{2em}} \nicefrac{X_1 \coprod X_2 }{\sim}=X
\]
mit $v_1 \sim f(v_1)$ für alle $v_1 \in E_1$ mit $p_1(v_1) \in X_0$ ein $K$-Vektorbündel über $X$.
\minisec{Beweis}
\emph{Übung!}
% subsection 37 (end)

\subsection{Lemma} % (fold)
\label{sub:38}
Sei $(X,A)$ ein Ko-Raumpaar. Sei $p \colon E \to A$ ein $K$-Vektorbündel und $\iota \colon A \hookrightarrow X$ die Inklusion. Dann existieren Vektorbündel $q \colon F \to X$
und $q' \colon F' \to A$, so dass $\iota^* F \cong E \oplus F'$.
\minisec{Beweisidee}
Finde $F' \xrightarrow{q'} A $, sodass $E \oplus F'$ trivialisierbar ist. Setze $F = X \times K^N$.
\[
	\begin{tikzcd}[row sep=0.1em]
		E\big|_{U_i} \rar & U_i \times K^n \\
		E \rar & X \times \Underbrace{K^n \oplus \ldots \oplus K^n}{\text{$k$ Kopien}}
	\end{tikzcd}
\]
% subsection 38 (end)

\subsection{Theorem} % (fold)
\label{sub:39}
\[
	\begin{tikzcd}
		X_0 \rar[hook,"i_1"] \dar["i_2"]& X_1 \dar["j_1"] \\
		X_2 \rar[hook,"j_2"] & X
	\end{tikzcd}
\]
Für das obige Pushout-Diagramm von kompakten Hausdorffräumen gilt
\begin{enumerate}[a)]
	\item Ist $i_1$ eine Kofaserung, dann auch $j_2$.
	\item Die Abbildungen $j_1^* \colon K^{-n}(X,X_2) \to K^{-n}(X_1,X_0)$ für $n \ge 0$ ist ein Isomorphismus.
\end{enumerate}
\minisec{Beweis}
\begin{enumerate}[a)]
	\item Da $i_1$ eine Kofaserung ist, gibt es die folgende Homotopieerweiterungseigenschaft
	\[
		\begin{tikzcd}
			X_1 \times \set{0} \cup X_0 \times I \rar \dar[hook] & Z \\
			X_1 \times I \urar[dashed,"\exists"'] 
		\end{tikzcd}
	\]
	Sei jetzt $X \times \set{0} \cup X_2 \times I \xrightarrow{f} Z$ gegeben ($f$ stetig). Wir haben 
	\[
		\begin{tikzcd}[column sep=huge]
			X_1 \times \set{0} \cup X_0 \times I \rar["{j_1 \cup (j_2 \times \id)}"] \dar[hook] & X \times \set{0} \cup X_2 \times I \rar["f"] & Z \\
			X_1 \times I \ar[ur, dashed,"h_1"'] & & 
		\end{tikzcd}
	\]
	Betrachte das Pushout-Diagramm
	\[
		\begin{tikzcd}
			X_0 \times I \rar[hook,"i_1 \times \id"] \dar["i_2 \times \id"] & X_1 \times I \dar["j_1 \times \id"] \ar[ddr, bend left,"h_1"] \\
			X_2 \times I \rar[hook] \ar[drr, bend right,"f|_{X_2 \times I}"] & X \times I \drar[dashed,"h"] \\
			& & Z
		\end{tikzcd}
	\]
	Also existiert $h$, sodass
	\[
		\begin{tikzcd}
			X \times \set{0} \cup X_2 \times I \rar["f"] \dar[hook]& Z \\
			X \times I \urar["h"'] 
		\end{tikzcd}
	\]
	Damit ist Teil a) gezeigt.
	\item Wir zeigen die Behauptung zunächst für $n=0$:
	\begin{description}[wide=0cm]
		\item[Injektivität:] Sei $(E_0, f,E_1)$ ein Tripel über $(X,X_2)$ mit $j_1^* \benbrace*{E_0,f,E_1}_\st =0$. Also existiert ein Vektorbündel $q \colon F \to X_1$, sodass
		sich der Isomorphismus
		\[
			(j_1^* E_0 \oplus F)\big|_{X_0} \xrightarrow{j_1^* f \oplus \id} (j_1^* E_1 \oplus F) \big|_{X_0}
		\]
		zu einem Isomorphismus über $X_1$ fortsetzen lässt. Nach Lemma \ref{sub:38} existieren Vektorbündel $q' \colon F' \to X$ und $\tilde{F} \to X_1$, sodass 
		$j_1^* F' \cong F \oplus \tilde{F}$ gilt. Es gilt
		\[
			\benbrace*{j_1^* E_0, j_1^* f , j_1^*E_1}_\st = \benbrace[\big]{j_1^* (E_0 \oplus F'), \Underbrace{j_1^*(f \oplus \id_{F'})}{\mathclap{\text{lässt sich zu Iso über 
			$X_1$ ausdehnen}}}, j_1^*(E_1 \oplus F')} 
		\]
		Wir können daher ohne Beschränkung der Allgemeinheit annehmen, dass sich $j_1^* f$ zu einem Isomorphismus über $X_1$ fortsetzen lässt. Das Vektorbündel $E_0$ ist 
		isomorph zur Verklebung von $j_1^* E_0 \to X_1$ und $j_2^* E_0 \to X_2$ mittels der Identität $i_2^* j_2^* E_0 \xrightarrow{\cong} i_1^* j_1^* E_0$ wie in Lemma 
		\ref{sub:37} beschrieben. Ähnliches gilt für $E_1 \to X$. Wir haben also Isomorphismen 
		\[
			j_1^* E_0 \xrightarrow[\cong]{\minwidthbox{j_1^* f }{2em}} j_1^* E_1 \quad \text{ und } \quad j_2^* E_0 = E_0\big|_{X_2} \xrightarrow[\cong]{\minwidthbox{f}{2em}}  E_1\big|_{X_2} = j_2^* E_1
		\]
		die kompatibel mit der Äquivalenzrelation in der Verklebung sind. Sie induzieren daher einen Isomorphismus $E_0 \xrightarrow{\cong} E_1$ über $X_1$, der $f$ fortsetzt.
		Es folgt $\benbrace*{E_0,f,E_1}_\st=0$.
		\item[Surjektivität:] Sei $(E_0,f,E_1)$ ein Tripel über $(X_1,X_0)$. Nach Lemma \ref{sub:38} existieren Vektorbündel $q \colon F \to X_1$ und $p_1' \colon E_1' \to X$
		mit einem Isomorphismus $\psi \colon E_1 \oplus  F \xrightarrow{\cong} j_1^* E_1'$. Sei $\alpha \colon (E_0 \oplus  F)\big|_{X_0} \to j_1^*E_1'\big|_{X_0}$ gegeben 
		durch
		\[
			(E_0 \oplus F) \big|_{X_0} \xrightarrow[\cong]{\minwidthbox{f \oplus  \id}{2em}} (E_1 \oplus F) \big|_{X_0} \xrightarrow[\cong]{\minwidthbox{\psi}{2em}} 
		\]
		Definiere nun $p_0' \colon E_0' \to X$ durch Verkleben von $E_0 \oplus F \to X_1$ und $E_1' \big|_{X_2} \to X_2$ entlang $\alpha$ wie in Lemma \ref{sub:37}. Sei 
		$f' \colon E_0'\big|_{X_2} \to E_1'\big|_{X_2}$ gegeben durch die Identität. Jetzt gilt
		\begin{align*}
			j_1^* \benbrace[\big]{E_0',f',E_1'}_\st = \benbrace[\big]{E_0 \oplus F, \psi \circ (f \oplus \id), j_1^* E_1'}_\st 
			&= \benbrace[\big]{E_0 \oplus F, f \oplus  \id, E_1 \oplus F}_\st \\
			&= \benbrace[\big]{E_0,f,E_1}_\st + \Underbrace{\benbrace[\big]{F,\id,F}_\st }{=0}
		\end{align*}
	\end{description} 
	Für den Fall $n>0$ lässt sich der Beweis auf den Fall $n=0$ zurückführen, denn 
	\[
		(X_1,X_0) \times (I^n,\partial I^n) = \enbrace*{X_1 \times I^n, X_1 \times \partial I^n \cup X_0 \times I^n},
	\]
	$X_0 \times I^n \cup X_1 \times \partial I^n \hookrightarrow X_1 \times I^n$ ist eine Kofaserung und das folgende Diagramm ist ein Pushout-Diagramm
	\[
		\begin{tikzcd}
			X_0 \times I^n \cup X_1 \times \partial I^n \rar[hook] \dar & X_1 \times I^n \dar \\
			X_2 \times I^n \cup X \times \partial I^n \rar[hook] & X \times I^n
		\end{tikzcd}
	\]
	Damit ist das Theorem bewiesen. \bewende
\end{enumerate}
% subsection 39 (end)

\subsection{Korollar} % (fold)
\label{sub:310}
Sie $(X,A)$ ein Ko-Raumpaar und $U \subseteq A$, sodass $\overline{U} \subset \mathring{A}$, dann ist $K^{-n}(X,A) \to K^{-n}(X \setminus U, A \setminus U)$ ein Isomorphismus.
\minisec{Beweis}
Setze $X_0 = A \setminus U$, $X_1 = X \setminus U$ und $X_2=A$ in Theorem \ref{sub:39}. Dann ist der Pushout gleich $X$. \bewende
% subsection 310 (end)

\section*{Produkte in $K$-Theorie und Kettenkomplexe von Vektorbündeln}
Klassen in $K^0(X)$ sind $[E_1,E_0]$ \enquote{$=$} $E_1 -E_0$, die Produkte könnten wir wie folgt definieren: 
\[
	[E_1,E_0] \cdot [E_1',E_0'] = \benbrace[\big]{(E_1\otimes E_1') \oplus (E_0 \otimes E_0'), (E_1 \otimes E_0')\oplus (E_0 \otimes E_1')} 
\]
Damit kann man aber nur schwer rechnen. Stattdessen wollen wir ein Produkt über Kettenkomplexe definieren.
\[
	\benbrace*{E_1,f,E_0}  \leadsto \begin{tikzcd}
		\cdots \rar & 0 \rar & E_1 \rar["\tilde{f}"] & E_0 \rar & 0 
	\end{tikzcd}
\]
wobei $\tilde{f}$ die Fortsetzung von $f$ über $X$ ist

\subsection*{Definition}
Ein Kettenkomplex von $K$-Vektorbündeln über einem Ko-Raumpaar $(X,A)$ ist ein Paar $(C_n,c_n)_{n \in \mathds{N}}$, wobei die $C_i$ endlichdimensionale Vektorbündel sind und 
die $c_i \colon C_n \to C_{n-1}$ Morphismen von $K$-Vektorbündeln mit folgenden Eigenschaften sind:
\begin{itemize}
	\item $\exists N \ge 0$ mit $C_n=0$ für $n >N$,
	\item $c_{n-1} \circ c_n =0$,
	\item für alle $x \in A$ ist der Kettenkomplex von Vektorräumen
	\[
		\begin{tikzcd}
			0 \rar & (C_n)\big|_{X} \rar["c_n|_X"] & C_{n-1}\big|_X \rar & \cdots \rar & C_1\big|_X \rar["c_1|_X"] & C_0\big|_X \rar & 0
		\end{tikzcd}
	\]
	exakt.
\end{itemize}
Eine Kettenabbildung zwischen $(C_*,c_*)$ und $(D_*,d_*)$ über $(X,A)$ ist eine Folge von Morphismen $f_k \colon C_k \to D_k$, so dass folgendes Diagramm kommutiert
\[
	\begin{tikzcd}
		\cdots \rar & C_k \rar["c_k"] \dar["f_k"] & C_{k-1} \rar \dar["f_{k-1}"]& \cdots \\
		\cdots \rar & D_k \rar["d_k"] & D_{k-1} \rar & \cdots
	\end{tikzcd}
\]
Zwei Kettenkomplexe $(C_*^i, c_*^i)$ über $(X,A)$ für $i \in \set{0,1}$ heißen \Index{homotop}, falls es einen Kettenkomplex $(C_*,c_*)$ über $(X\times I, A \times I)$ gibt,
sodass $(C_*,c_*)\big|_{X \times \set{i} } \cong (C_*^i,c_*^i)$. Wir schreiben $(C_*^0,c_*^0) \simeq (C_+^1, c_*^1)$.

\todo[inline]{Ab hier fehlen noch diverse Diagramme}
\subsection*{Definition}
Sei $K(X,A)^c$ der Quotient der freien abelschen Gruppe erzeugt von den Isomorphie-Klassen von Kettenkomplexen von $\mathds{C}$-Vektorbündeln über $(X,A)$ bezüglich der 
folgenden Relationen:
\begin{enumerate}[a)]
	\item Wenn $C_* \simeq D_*$ gilt, dann folgt $[C_*] = [D_*]$
	\item Wenn $\begin{tikzcd}[cramped,sep=small]
		0 \rar & C_* \rar & D_* \rar & E_* \rar & 0
	\end{tikzcd}$ eine kurze exakte Folge von Kettenkomplexen ist (faserweise exakt), dann gilt
	\[
		[C_*] - [D_*] + [E_*] = 0
	\]
	\item \(
		\big[ \begin{tikzcd}[cramped,sep=small]
			\cdots \rar & 0 \rar & E \rar["\id"] & E \rar & 0 \rar & \cdots
		\end{tikzcd}\big]
	\) \todo{Hier fehlt noch eine Aussage \ldots}
\end{enumerate}
\minisec{Bemerkung}
Aus b) folgt $[C_*] = [D_*]$ für isomorphe Kettenkomplexe und 
\[
	\benbrace*{C_* \oplus  D_*, c_* \oplus d_*} = [C_*,c_*] + [D_*,d_*] 
\]
Wir haben einen Gruppenhomomorphismus $a \colon K^0(X,A) \to K(X,A)^c$, 
% $[E_1,f,E_0]_\st \mapsto  \benbrace*{\begin{tikzcd}[cramped,sep=small] \cdots \rar & 0 \rar & E_1 \rar["\tilde{f}"] & E_0 \rar & 0 \end{tikzcd}} $,
wobei $\tilde{f}$ eine Fortsetzung von $f$ auf $X$ ist. Dies ist unabhängig von der Wahl der Fortsetzung. Seien $\tilde{f}_0, \tilde{f}_1$. Dann sei 
$h_t = t \tilde{f}_1 + (1-t) \tilde{f}_0$. Es gilt $h_t|_A =f$ und 
\[
	\begin{tikzcd}[sep=5.5em]
		\cdots \rar & 0  \rar & E_1 \times I \rar["{(v,t) \mapsto (h_t(v),t)}"] & E_0 \times I \rar & 0
	\end{tikzcd}
\]
ist ein Kettenkomplex über $(X \times I, A \times I)$ und somit eine Homotopie zwischen 
$\begin{tikzcd}[cramped,sep=small]0 \rar & E_1 \rar["\tilde{f}_0"] & E_0 \rar & 0\end{tikzcd}$ und 
% $\begin{tikzcd}[cramped,sep=small]
% 	0 \rar & E_1 \rar["\tilde{f}_1"] & E_0 \rar 0
% \end{tikzcd}$

\subsection{Definition 3.1.2} % (fold)
\label{sub:def_312}
Eine \Index{Kettenkontraktion} über $A$ ist eine Familie von Bündelabbildungen $\gamma_n \colon C_n|_A \to C_{n+1}|_{A}$, sodass 
$\gamma_{n-1} \circ (c_n|_A) + (c_{n+1}|_A) \circ \gamma_n = \id_{C_n|_A}$ 
% \[
% 	\begin{tikzcd}
% 		\cdots \rar & C_{n+1}\big|_{A} \rar["c_{n+1}"] & C_n\big|_A  \dar[dashed,"\id"] \dlar[$\gamma_n$] \rar["c_n"] & C_{n-1}\big|_A \rar \dlar["\gamma_{n-1}"] \cdots \\
% 		\cdots \rar & C_{n+1}\big|_A \rar[c_{n+1}] & C_n\big|_A \rar & C_{n-1}\big|_A \rar \cdots
% 	\end{tikzcd}
% \]
% subsection def_312 (end)

\subsection{Lemma 3.1.3} % (fold)
\label{sub:lemma_3_1_3}
Sei $(C_*,c_*)$ ein Kettenkomplex über dem Ko-Raumpaar $(X,A)$. Dann existiert eine Kettenkontraktion über $A$.
\minisec{Beweis}
\emph{Übung.}
% subsection lemma_3_1_3 (end)

\subsection{Lemma 3.1.4} % (fold)
\label{sub:lemma_3_1_4}
Die Abbildung $b \colon K(X,A)^c \to K^0(X,A)$ mit $b([C_+,c_*]) = \enbrace*{C_{\mathrm{odd}}, c_n|_A + \gamma_*, C_{\mathrm{even}}}$ mit 
$C_{\mathrm{odd}} = \bigoplus_{k \in \mathds{N}_0} C_{2k+1}$ und $C_{\mathrm{even}} = \bigoplus_{k \in \mathds{N}_0} C_{2k}$, wobei $\gamma_*$ eine Kettenkontraktion über 
$A$ ist, ist wohldefiniert.
\minisec{Beweis}
Sei $\gamma_n' \colon C_n|_A \to C_{n+1}|_A$ eine weitere Kettenkontraktion über $A$. Betrachte 
$\delta_n = (\gamma_{n+1}' - \gamma_{n+1}) \circ \gamma_n \colon C_n|_A \to C_{n+2}|_A$ und
\begin{align*}
	&f \colon C_{\mathrm{even}}\big|_A \xrightarrow{\id + \delta_*} C_{\mathrm{even}}\big|_A \qquad  \text{ und} \\
	& g  \colon C_{\mathrm{odd}}\big|_A \xrightarrow{c_*|_A + \gamma_*} C_{\mathrm{even}}\big|_A \xrightarrow{\minwidthbox{f}{2em}} C_{\mathrm{even}}\big|_A \xrightarrow{c_*|_A + \gamma_*'}
	C_{\mathrm{odd}}\big|_A    
\end{align*}
Behauptung: $g$ ist ein Isomorphismus.
\missingfigure{Diagramm mit Vektorschreibweise}
Der Teil von $g$, der $C_{2n+1}|_A \to C_{2n-1}|_A$ abbildet, ist gegeben durch $c_{2n+1}|_A \circ c_{2n}|_A=0$. Der Teil von $g$, der $C_{2n+1}|_A \to C_{2n+1}|_A$ abbildet 
ist gegeben durch die Identität, also ist $g$ eine \enquote{Dreiecksmatrix mit Identitäten auf der Diagonalen} und somit ein Isomorphismus. Insgesamt folgt, dass
$c_*|_A + \gamma_*$ injektiv ist und $c_*|_A + \gamma_*'$ surjektiv. Aus Symmetriegründen sind beide Isomorphismen.

Betrachte jetzt $h_t = (\id + t \cdot \delta) \circ (c_+|_A + \gamma_*) \colon C_{\mathrm{odd}}|_A \to C_{\mathrm{even}}|_A$. Dies ist eine Homotopie durch Isomorphismen. 
Aufgrund der Homotopieinvarianz von $K^0(X,A)$ gilt:
\[
	\benbrace*{C_{\mathrm{odd}}, (c_*|_A + \gamma_*), C_\mathrm{even}}_\st  = \benbrace*{C_{\mathrm{odd}}, f \circ (c_*|_A + \gamma_*), C_\mathrm{even}}_\st
\] 
Auch die Kettenabbildung $g$ ist von der Form $g = \id + \varepsilon$, wobei $\varepsilon$ der Teil oberhalb der Diagonalen ist. Mit dem gleichen Trick wie eben erhalten wir
mit $\tilde{h}_t= (c_*|_A + \gamma_*') ^{-1} \circ (\id + t \varepsilon)$
\[
	\benbrace[\big]{C_\mathrm{odd}, \Underbrace{(c_*|_A + \gamma_*') ^{-1}\circ g}{=f \circ (c_*|_A + \gamma_*)}, C_\mathrm{even}}_\st = \benbrace*{C_\mathrm{odd}, (c_*|_A + \gamma_*') ^{-1}, C_\mathrm{even}}_\st  
\]
Also insbesondere für $\gamma_*' = \gamma_*$:
\[
	\benbrace*{C_\mathrm{odd}, (c_*|_A + \gamma_*), C_\mathrm{even}}_\st = \benbrace*{C_\mathrm{odd}, (c_*|_A + \gamma_*) ^{-1}, C_\mathrm{even}}  
\]
Aus den letzten drei Formeln\todo{sternchen oder sowas einfügen} folgt 
\[
	\benbrace*{C_\mathrm{odd}, (c_*|_A + \gamma_*), C_\mathrm{even}} = \benbrace*{C_\mathrm{odd}, (c_*|_A + \gamma_*')^{-1}, C_\mathrm{even}}  
	= \benbrace*{C_\mathrm{odd}, (c_*|_A + \gamma_*'), C_\mathrm{even}} 
\]
Kompatibilität mit den Relationen:
\begin{itemize}
	\item $C_* \simeq D_* \Rightarrow b([C_*]) = b([D_*])$ aus der Homotopieinvarianz von $K^0(X,A)$.
	\item $b \enbrace*{\benbrace*{\cdot \to 0 \to E \xrightarrow{\id} E \to 0 \to \cdots } } = \benbrace*{E,\id,E}_\st = 0$
	\item Sei $0 \to C_* \xrightarrow{i_*} D_* \xrightarrow{p_*} E_0 \to 0$ eine kurze exakte Sequenz von Kettenkomplexen. Jede kurze exakte Sequenz von Vektorbündeln 
	spaltet (Übung!).
	Also existieren Morphismen $t_n \colon E_n \to D_n$ mit $p_n \circ t_n = \id_{E_n}$. Problem: $t_n$ ist keine Kettenabbildung! Wähle eine Kettenkontraktion $\varepsilon$
	über $A$ für $E_*$. Setze
	\[
		s_n = d_{n+1}|_A \circ t_{n+1}|_A \circ \varepsilon_n + e_n|_A \circ \varepsilon_{n-1} \circ e_n |_A
	\]
	Dies ist ein Spalt von $p_n$ über $A$. $i_n +s_n \colon C_n|_A \oplus E_n|_A \to D_n|_A$ ist homotop durch Isomorphismen zu $i_n + t_n$ via
	\[
		h_r = r \enbrace*{d_{n+1}\big|_A \circ t_{n+1}\big|_A - t_n\big|_A \circ e_{n+1}\big|_A} \circ \varepsilon_n + (t_n - i_n) \big|_A
	\]
	Da $i_n,t_n$ auf ganz $X$ definiert sind, lässt sich $i_n+ s_n$ zu einem Isomorphismus $\varphi_n$ über $X$ fortsetzen. Setze 
	$\delta_n = (i_n +s_n) \circ (\gamma_n \oplus \varepsilon_n) \circ (i_n +s_n)^{-1}$. Dies ist eine Kettenkontraktion von $D_*$ über $A$. Es folgt
	\begin{align*}
		\benbrace[\big]{C_\mathrm{odd}\oplus E_\mathrm{odd}, (c_*|_A + \gamma_*) \oplus (e_*|_A + \varepsilon_*), C_\mathrm{even} \oplus E_{\mathrm{even}}}_\st
		= \benbrace[\big]{D_\mathrm{odd}, d_*|_A + \delta_*, D_\mathrm{even}}_\st \bewende  
	\end{align*}
\end{itemize}
% subsection lemma_3_1_4 (end)

\subsection{Lemma 3.1.5} % (fold)
\label{sub:lemma_3_1_5}
Seien $(E_*,e_*)$ und $(E_*,e_*')$ Kettenkomplexe über $(X,A)$ mit $e_i|_A = e_i'|_A$. Dann gilt $\benbrace*{E_*,e_*} = \benbrace*{E_*,e_*'} \in K(X,A)^c$.
% subsection lemma_3_1_5 (end)

\subsection{Lemma 3.1.6} % (fold)
\label{sub:lemma_3_1_6}
Die Abbildungen $a$ und $b$ sind invers zueinander.
\minisec{Beweis}
$b \circ a= \id_{K^0(X,A)}$ ist klar. Es bleibt zu zeigen, dass $a$ surjektiv ist $\iff$ Jeder Kettenkomplex der Länge $l \ge 2$ ist äquivalent zu einem Kettenkomplex der
Länge $l=1$. Es gilt
\begin{align*}
	&\benbrace*{0 \to C_n \xrightarrow{c_n} C_{n-1} \to C_{n-2} \to \ldots \to C_0 \to 0 } \\
	= &\benbrace*{0 \to C_n \xrightarrow{c_n \oplus 0} C_{n-1} \oplus C_n \xrightarrow{c_{n-1} \oplus \id} C_{n-2} \oplus C_n \to C_{n-3} \to \ldots \to C_0 \to 0 } 
\end{align*}
\begin{itemize}
	\item $(c_n \oplus 0)|_A$ ist ein Monomorphismus und homotop zu $(0 \oplus \id)|_A$ durch Monomorphismen. Jetzt lässt sich $(0 \oplus \id)|_A$ zu einem Monomorphismus 
	über $X$ ausdehnen. Eine Variante von Lemma \ref{sub:33} zeigt, dass sich dann auch $(c_n \oplus 0)|_A$ zu einem Monomorphismus $\tau \colon c_N \to C_{n-1} \oplus C_n$
	auf $X$ ausdehnen lässt.
	\[
		\Rightarrow C_{n-1} \oplus C_n \cong \tau (C_n) \oplus  Q
	\]
	Der Kettenkomplex 
	\[
		0 \to C_n \xrightarrow{\tau} \tau(C_n) \oplus Q \xrightarrow{(c_{n-1} \oplus \id)|_Q \circ \pr_Q} C_{n-2} \oplus C_n \to \ldots \to C_0 \to 0  
	\]
	stimmt über $A$ mit dem obigen überein, liefert also dieselbe Klasse. Aber
	\begin{align*}
		\Underbrace{\benbrace*{0\to C_n \xrightarrow{\tau} \tau(C_n) \to 0  \to \cdots }}{= \benbrace*{0 \to C_n \xrightarrow{\id} C_n \to 0 \to \ldots \to 0 }=0 } + 
		\Underbrace{\benbrace*{0 \to Q \xrightarrow{C_{n-1}'} C_{n-2} \oplus C_n \to \ldots \to C_0 \to 0 } }{\text{Länge um 1 verringert}} \bewende
	\end{align*}
\end{itemize}
% subsection lemma_3_1_6 (end)

Sei $(C_.,c_.)$ ein Kettenkomplex über dem Ko-Raumpaar $(X,A)$ und $(C_.',c_.')$ ein weiterer über $(X',A')$. Das \Index{äußere Tensorprodukt} von zwei Vektorbündeln $E \to X$,
$F \to X'$ ist definiert durch 
\[
	E \boxtimes F := \pi_X^* E \otimes \pi_{X'}^* F
\]
wobei $\pi_X \colon X \times X' \to X$, $\pi_{X'} \colon X \times X' \to X'$ die Projektionen sind. Es gilt also $(E \boxtimes F)_{x,x'} = E_x \times F_{x'}$.
Das \Index{Tensorprodukt} von $(C_.,c_.)$ und $(C_.',c_.')$ ist definiert durch 
\[
	(C \boxtimes C')_n := \bigoplus_{i+j=n} C_i \boxtimes C_j'
\]
mit dem Differential $c_n''$, das auf $C_i \boxtimes C_j'$ gegeben ist durch 
\[
	c_n'' \big|_{C_i \boxtimes C_j'} = c_i \boxtimes \id + (-1)^i \id{} \boxtimes c_j'
\]
Damit ist $\enbrace*{(C \boxtimes C')_., c_.''} $ ein Kettenkomplex über $(X,A) \times (X',A')$. Somit erhalten wir 
\[	  
	\begin{tikzcd}
		K^{-n}(X,A) \times K^{-m}(X',A') \dar["a \times a","\cong"'] \rar["*"] & K^{-(n+m) \enbrace*{(X,A) \times (X',A')} } \dar["a","\cong"'] \\
		K \enbrace*{(X,A) \times (I, \partial I)^n}^c \times K \enbrace*{(X',A') \times (I, \partial I)^m}^c \rar &
		K^{-(n+m)} \enbrace*{(X,A) \times (X',A') \times (I,\partial I)^{n+m}}    
	\end{tikzcd}
\]
Für $X=X'$ und $A \subset X, B \subset X$ sei 
$\Delta \colon (X,A \cup B) \to (X,A) \times (X,B)= (X\times X, X \times B \cup A \times X) \supset (A \cup B) \times (A \cup B)$. Dies induziert 
$K^{-n(n+m)} \enbrace*{(X,A) \times (X,B)} \to K^{-(n+m)}(X, A \cup B) $. Wir definieren
\[
	\cdot := \Delta^* \circ * \colon K^{-n}(X,A) \times K^{-m}(X,B) \xrightarrow{*} K^{-(n+m)} \enbrace*{(X,A) \times (X,B)}\xrightarrow{\Delta^*} K^{-(n+m)}(X,A \cup B)   
\]
\begin{itemize}
	\item Insbesondere ist $K^*(X) = \bigoplus_{n=0}^\infty K^{-n}(X)$ ein graduiert kommutativer Ring mit $1$, das heißt es gilt für $x \in K^i(X)$ und $y \in K^j(X)$
	\[
		x \cdot y = (-1)^{i \cdot j} y \cdot x
	\]
	\item Für $f \colon X \to Y$ erhalten wir einen Ringhomomorphismus $f^* \colon K^*(Y) \to K^*(X)$.
	\item Für $A \subset B \subset X$ ist $K^*(B,A)$ ein Modul über $K^*(X)$.
\end{itemize}

\subsection*{Lemma}
Für $a,a' \in K^0(X)$, $b,b' \in K^0(Y)$ ist 
\[
	(a \cdot a') * (b \cdot b') = (a * b) \cdot (a' * b')
\]
\minisec{Beweis}
\emph{Richtiges Diagramm hinschreiben.}

\newpage
\section*{3.2 Bott-Periodizität -- eine Beweisskizze}
Welche $K$-Gruppen können wir berechnen? 
\begin{itemize}
	\item $K^0(\pt) = \Gr(\mathds{N}) \cong \mathds{Z} $
	\item $K^{-1}(\pt) = K^0(I,\partial I) = 0$
	\item Tripel $(E_1,f,E_0)$ über $(I,\partial I)$ mit $f \colon E_1|_{\partial I} \xrightarrow{\cong} E_0|_{\partial I}$. Ohne Beschränkung der Allgemeinheit sei 
	$E_i \cong \mathds{C}^n \times I$. Da $\GL_n(\mathds{C})$ wegzusammenhängend ist, gibt es eine Erweiterung von $f$ zu einem Isomorphismus über $I$. 
	\[
		f_0 \colon \mathds{C}^n \times \set{0} \xrightarrow{\cong} \mathds{C}^n \times \set{0} \\
		f_1 \colon \mathds{C}^n \times \set{1} \xrightarrow{\cong} \mathds{C}^n \times \set{1}     
	\]
	$f$ entspricht also zwei Elementen von $\GL_n(\mathds{C})$. $g \colon \mathds{C}^n \times I \xrightarrow{\cong}  \mathds{C}^n \times I \leadsto I \to \GL_n(\mathds{C})$ 
\end{itemize}
Übung $KO ^{-1}(\pt)$ ?

Was ist $K^0(S^1)$? 
\[
	\begin{tikzcd}[row sep=0.4em] 
		& \benbrace*{S^1 \times \mathds{C},0}_\st \rar[mapsto] & \benbrace*{\mathds{C},0}_\st = 1  \\
		0=K^0(I,\partial I) \cong K^0(S^1,\pt) \rar & K^0(S^1) \rar & K^0(\pt) \cong \mathds{Z}
	\end{tikzcd}
\]
Aus der unteren Zeile folgt $ K^0(S^1)$ ist Untergruppe von $\mathds{Z}$. Aus der oberen folgt dann weiter $K^0(S^1)\cong \mathds{Z}$.

Gibt es eine nichttriviale Klasse in $K^0(S^2)$? Wir haben einen Homöomorphismus 
\[
	S^2 \cong \mathds{C}P^1 = \set[V \subset \mathds{C}^2]{V \text{ ist eindimensionaler Untervektorraum}} 
\]
Betrachte $H= \set[(v,V) \in \mathds{C}^2 \times \mathds{C}P^1]{v \in V} $. Dies ist ein eindimensionales Vektorbündel über $\mathds{C}P^1$. Was ist 
$[H,0]_\st \in K^0(\mathds{C}P^1)$? Welche Gestalt haben Vektorbündel über $X \times S^2$? Sei $S^2= D_+ \cup D_-$ die Zerlegung in die \enquote{obere} und die 
\enquote{untere} Halbsphäre. Sei $E \to X \times S^2$ ein Vektorbündel. Sei $z_0 \in S^1= D_+ \cap D_- \subset S^2$ ein Basispunkt von $S^2$. Sei 
$F= \iota^* E$ mit $\iota \colon X \to X \times S^2$ gegeben durch $x \mapsto(x,z_0)$. Wegen $D_+ \cong D_- \simeq \set{\pt}$ ist $\pr_X^\pm \colon X \times D_{\pm} \to X$
eine Homotopieäquivalenz mit Homotopieinversem $\iota_{\pm} \colon X \to X \times D_\pm$, $x \mapsto (x,z_0)$. Daher ist $E|_{X \times D_\pm} \cong \iota_\pm^*F$.

Dann ist $E$ isomorph zur Verklebung von $\iota_+^*F$ und $\iota_-^* F$ entlang eines Isomorphismus, der gegeben ist durch eine stetige Abbildung $S^1 \to \textsc{Iso}(F)$, 
wobei $\textsc{Iso}(F)$ der Raum aller Vektorbündelisomorphismen $F \to F$ ist. Wir schreiben $E \cong \set{F,f}$. Die Isomorphieklasse von $\set{F,f}$ hängt nur noch der 
Homotopieklasse von $f$ ab.
\missingfigure{Diagramm}

\minisec{Beispiel}
Übergangsabbildung von $H$: Identifiziere $\mathds{C}P^1$ mit $\mathds{C} \cup \set{\infty} \cong S^2 $ mit Hilfe der Abbildung 
$\mathds{C} \cup \set{\infty} \to \mathds{C}P^1 $, definiert duch $z \mapsto \benbrace*{z:1}, \infty \mapsto \benbrace*{1:0}$. Die Umkehrabbildung ist gegeben durch
\[
	\benbrace*{z_1,z_2} \mapsto \begin{cases}
		z_1/z_2, &\text{ falls }z_2 \not=0\\
		\infty , & \text{ sonst}
	\end{cases}
\]
Nun ist
\begin{align*}
	H &= \set[(z,v) \in (\mathds{C} \cup \set{\infty}) \times \mathds{C}^2]{v= \lambda(z,1), \text{ falls } z \not= \infty \text{ oder } v=(\lambda,0) \text{ falls } 
	z=\infty} \\
	D_+ &= \set[z \in (\mathds{C} \cup \set{\infty})]{\abs*{z} \le 1 } \\
	D_- &= \set[z \in (\mathds{C} \cup \set{\infty})]{ \abs*{z} \ge 1 \text{ oder } z= \infty }  
\end{align*}
Trivialisierungen:
\[
	\begin{tikzcd}
		H\big|_{D_+} \rar["\cong"] & D_+ \times \mathds{C} \\
		(z, (v_1,v_2)) \rar[mapsto] & (z,v_2) \\
		(z,(wz,w)) & (z,w) \lar[mapsto]
	\end{tikzcd}
	, 
	\begin{tikzcd}
		H\big|_{D_-} \rar["\cong"] & D_- \times \mathds{C} \\
		(z, (v_1,v_2)) \rar[mapsto] & (z,v_1) \\
		(z,(w,w/z)) & (z,w) \lar[mapsto]
	\end{tikzcd}
\]
Wir setzen $F=S^1  \times \mathds{C}$ und $f \colon S^1 \to \textsc{Iso}(S^2 \times \mathds{C})$, $z \mapsto \enbrace*{w \mapsto wz}$, in jeder Faser Multiplikation 
mit $z$. Es ist nun $H= \set{\mathds{C}, w \mapsto w \cdot z} = \set{\mathds{C},z} $ und $H^* \cong \set{\mathds{C},z ^{-1}}$.

Weitere Beispiele
\[
	H^{\otimes K} = \set{\mathds{C}, w \mapsto w \cdot z^k} \quad , \quad H \oplus H = \set{\mathds{C} \oplus \mathds{C}, z \mapsto \begin{pmatrix}
		z & 0 \\ 0 & z \end{pmatrix}} 
\]
Da $\GL_2(\mathds{C})$ wegzusammenhängend ist, gibt es einen Pfad zwischen $\begin{psmallmatrix}1 & 0 \\ 0 & 1 \end{psmallmatrix}$ und
$\begin{psmallmatrix} 0& 1 \\ -1 & 0 \end{psmallmatrix}$. Betrachte 
\todo[inline]{hier fehlt noch was}
% \[
% 	\begin{tikzcd}
% 		S^1 \times I \rar & \GL_n(\mathds{C})
% 		(z,t) \rar[mapsto] &\begin{pmatrix} z & 0 \\ 0 & 1 \end{pmatrix} \gamma_t \cdot
% 	\end{tikzcd}
% \]
Hieraus erhalten wir
\[
	H^{\otimes 2} \oplus \Underbrace{(\mathds{C} \times S^2)}{\cong \set{C \oplus \mathds{C}, \begin{pmatrix}
		z^2 & 0 \\ 0 & 1
	\end{pmatrix}}} \cong \Underbrace{H \oplus  H}{ \cong \set{\mathds{C} \oplus  \mathds{C}, \begin{pmatrix}
		z & 0 \\
		0 & z
	\end{pmatrix}}  }
\]
In $K$-Theorie
\[
	\benbrace*{H,0} \cdot \benbrace*{H,0} + \benbrace*{\mathds{C},0}    = \benbrace*{H^{\otimes 2} \oplus \mathds{C},0}_\st 
	= \benbrace*{H \oplus H,0}_\st = 2 \cdot \benbrace*{H,0} 
\]
in $K^0(S^2)$ gilt die Relation $\benbrace*{H,0}^2 - 2 \benbrace*{H,0} + \benbrace*{\mathds{C},0}= 0 $
\[
	(x-1)^2 = X^2 - 2x +1 = 0
\]

\subsection*{Theorem Bott '58}
\[
	\begin{tikzcd}
		\mu \colon K^1(X) \otimes \nicefrac{\mathds{Z}[H]}{(H-1)^2} \rar & K^0(X) \otimes  K^0(S^2) \rar & K^0(X \times S^2) \\
		\benbrace{E,0}_\st \otimes H^m \ar[rr,mapsto] & & \benbrace[\big]{\Underbrace{\pr_X E \otimes (\pr_{S^2}^* H)^{\otimes m}}{\set{E,z^m} },0}_\st
	\end{tikzcd}
\]
ist ein Ringisomorphismus . Insbesondere $K^0(S^2) \cong \nicefrac{\mathds{Z}[H]}{(H-1)^2}$
\minisec{Beweis}
\begin{enumerate}[1.]
	\item Satz: Jedes Vektorbündel $\set{F,f}$ über $X \times S^2$ ist isomorph zu $\set{F,l}$, wobei $l= \sum_{i=-N}^{N} a_i \cdot z^i$ und $a_i \colon F \to F$ ein 
	Morphismus von Vektorbündeln ist, das heißt die Übergangsfunktion ist gegeben durch ein Laurent-Polynom. Sind $l_0$ und $l_1$ homotope Übergangsfunktionen, dann sind sie 
	auch homotop durch Laurent-polynomiale Übergangsfunktionen.
	
	Idee hierzu: Wir können stetige Funktionen $g \colon X \times S^1 \to \mathds{C}$ durch Laurent-Polynome approximieren. 
	\(
		g(x,z) = \sum_{n \in \mathds{Z}} a_n(x) z^n
	\)
	für 
	\[
		a_n(x) = \frac{1}{2 \pi } \int_0^{2 \pi} g(x, e^{i \theta}) e^{- i \theta n} \mathd \theta 
	\]
	Genauso können wir Funktionen $\tilde{g} \colon X \times S^1 \to M_n \mathds{C}$ approximieren. Wähle ein Skalarprodukt in den Fasern von $F$ (hermitesche Metrik).
	Der Raum der Vektorbündelendomorphismen $\END(F \times S^1)$ ist normiert via
	\[
		\norm{\alpha} = \sup_{\substack{\norm{v}=1 \\ v \in F \times S^1 }} \norm{\alpha(v)} 
	\]
	$\Iso(F \times S^1) \subset \END(F \times S^1)$ ist offen. Fasse $f$ als Schnitt $f \colon X \times S^1 \to \Iso(F \times S^1) \subset \End(F \times S^1)$. Wähle 
	trivialisierende Überdeckung $U_i$ mit Partition der Eins $\set{\varphi_i}_i$
	\[
		f \big|_{U_i} \colon U_i \times S^1 \longrightarrow \End \enbrace*{F\big|_{U_i} \times S^1} \cong M_n \mathds{C} \times U_i \times S^1 
	\]
	Approximiere durch $l_i$ und setze $l= \sum_{i} \varphi_i l_i$. $l$ liegt beliebig nahe bei $f$ in Norm. Da $\Iso$ offen ist, folgt dass $f$ und $l$ homotop durch
	Isos. \bewende
	\item (Linearisierungstrick): Schreibe $\set{F,l}= \set{F, z^{-m} \cdot q}$ für ein \enquote{Polynom} $q$ und ein $m \in \mathds{N}$.
	
	Satz: Sei $q$ eine polynomiale Übergangsfunktion vom Grad $\le n$, dann gilt 
	\[
		\mengenA{\Underbrace{n \cdot F}{F \oplus \ldots \oplus F},1} \oplus \set{F,q} \cong \set{(n+1) \cdot F, L^n q} 
	\]
	wobei $L^n q$ eine Übergangsfunktion ist, die nur noch $z^0$ und $z^1$ enthält, das heißt $L^n q(x,z) = a(x) + b(x) \cdot z$.
	
	Idee hierzu: $q= \sum_{i=1}^{n} a_i z^i$. Übergansfunktion der linken Seite ist gegeben durch eine Matrix
	\[
		\begin{pmatrix}
			1 & & & 0 \\
			 & 1 & &  \\
			 & & \ddots &  \\
			 0 & & & q
		\end{pmatrix}
	\]
	Zeilen- und Spaltenoperationen bringen dies auf die Form
	\[
		\begin{pmatrix}
			1 & -z  & & &\\
			& 1 & -z & &  \\
			& & \ddots & \ddots & \\
			& & & 1 & -z \\
			a_n & a_{n-1} & \cdots & a_1 & a_0
		\end{pmatrix} = A
	\]
	Zeilen- und Spaltenoperationen erhalten die Determinante, also ist $A$ invertierbar für alle $(x,z) \in X \times S^1$. Damit ist $A$ eine gültige Übergangsfunktion.
	Es ist $A = a + b \cdot z= L^n q$. Jede Zeilen- und Spaltenoperation ist homotop zur Identität, daher ist $L^n q$ homotop zu $1 \oplus q$
	\[
		\Rightarrow \set{(n+1) \cdot F,1 \oplus q} \cong \set{(n+1)\cdot F, L^n q} \bewende  
	\]
	\item (Aufspaltung von $F$): Wir haben bereits: $\set{(n+1) \cdot F, q} \cong \set{(n+1) \cdot F, a+ b \cdot z}$. Betrachte nun die Homotopie 
	\[
		q_t(x,z)=  (1 + tz) \enbrace*{b(x) \frac{z+t}{1+ tz} + a(x)} = \enbrace*{b(x) + t a(x)} z + t b(x) + a(x) 
	\]
	Der mittlere Teil ist gleich $L^n q$ für $t=0$, invertierbar für $0 \le t <1$. Weiter ist $z \mapsto \frac{z+t}{1+ tz}$ ein Homöomorphismus $S^1 \to S^1$.
	Wähle nun $t_0 <1$, sodass $b(x) + t_0 a(x)$ invertierbar ist. Dies funktioniert, denn $b(x) + a(x)$ ist invertierbar. Damit erhalten wir 
	\[
		\set{(n+1) \cdot F, L^n q} \cong \set{(n+1) \cdot f, z+c}  
	\]
	mit $c(x)= \enbrace*{t_0 b(x) + a(x)}\enbrace*{b(x) + t_0 a(x)}^{-1}$, $b+ t_0 \cdot a \colon F \xrightarrow{\cong} F$. Insbesondere ist 
	$z+ c(x) \colon F_x \times \set{z}  \xrightarrow{\cong} F_x \times \set{z}$ ein Isomorphismus, also hat $z + c(x)$ niemals Eigenwert $0$, woraus folg, dass $c(x)$ keine
	Eigenwerte auf $S^1$ hat.
	
	Lemma: Sei $c \colon F \to F$ ein Morphismus von Vektorbündeln, so dass $c(x) \colon F_x \to F_x$ keine Eigenwerte auf $S^1$ hat. Dann existieren eindeutige 
	Untervektorbündel $F_+$ und $F_-$, sodass 
	\begin{enumerate}[a)]
		\item $F \cong F_+ \oplus F_-$
		\item $c(F_\pm) \subset F_\pm$
		\item Die Eigenwerte von $c|_{F_+}$ liegen außerhalb von $S^1$, von $c|_{F_-}$ innerhalb von $S^1$
	\end{enumerate}
	
	Korollar: $\set{F,  z+c} \cong \set{F_+,1} \oplus \set{F_-,z}$. Beweis: 
	\[
		\set{F, z+c} \cong \set{F_+, z+ c_+} \oplus \set{F_-, z+ c_-} 
	\]
	mit $c_\pm = c|_{F_\pm}$. Die Eigenwerte von $z+ c_+$ liegen außerhalb von $S^1$, $tz + c_+$ liefert eine Homotopie zu $\set{F_+, c_+} \cong \set{F_+,1}$.
	Die Eigenwerte von $c_-$ liegen innerhalb von $S^1$. Dann liefert $z + t c_-$ eine Homotopie durch Isomorphismen zu $z$, das heißt $\set{F_-, z + c_-} \cong\set{F_-,z}$.
	\bewende
\end{enumerate}
Setze alles zusammen für die Surjektivität von $\mu$. $\mu$ ist ein Gruppenhomomorphismus, also müssen wir nur zeigen, dass 
$K^0(X \times S^2)\ni \benbrace*{\set{F,f},0 }_\st $ im Bild liegt. Es gilt nun
\begin{align*}
	\benbrace*{\set{F,f},0 }_\st \stackrel{1}{=} \benbrace*{\set{F, z^{-m} q},0 }_\st &= \benbrace*{\set{F,q} \oplus H^{-m},0 }_\st \\
	&\stackrel{2}{=} \benbrace*{\set{(n+1) F, L^n q} \oplus H^{-m},0 }_\st -  \benbrace*{\set{n F,1} \otimes H^{-m}, 0 }_\st \\
	&\stackrel{3}{=} \benbrace*{ \set{\enbrace*{(n+1)F}_+,1} \otimes H{-m},0 }_\st + \benbrace*{ \set{\enbrace*{(n+1)F}_-,z }\otimes H{-m},0 }_\st \\
	&\quad - \benbrace*{\set{n F,1} \otimes H^{-m},0 }_\st   \\
	&=  \benbrace*{\pr_X^* \enbrace*{(n+1)F}_+ \otimes H^{-m},0 }_\st + \benbrace*{\pr_X^* \enbrace*{(n+1)F}_- \otimes H^{1-m},0 }_\st  \\
	&\quad - \benbrace*{\pr_x^*(n F) \otimes H^{-m},0}_\st 
\end{align*}
liegt im Bild von $\mu$.
Für Injektivität betrachte die Abbildung 
\[
	\begin{tikzcd}[row sep=0.0em]
		\nu \colon K^0(X \times S^2) \rar & K^0(X)\otimes \nicefrac{\mathds{Z}[H]}{(H-1)^2} \\
		\benbrace*{\set{F,f},0 }_\st \rar[mapsto] & \benbrace*{ \enbrace*{(n+1)F}_- ,0 }_\st \otimes (H-1) + \benbrace*{F ,0}_\st \otimes H^{-m} 
	\end{tikzcd}
\]
Prüfe, dass dies wohldefiniert ist! $\nicefrac{\mathds{Z}[H]}{(H-1)^2}$ wird erzeugt von $1, H ^{-1}$. Dann gilt
\begin{align*}
	\nu \circ \mu \enbrace*{ \benbrace*{F,0}_\st \otimes H^{-m}} = \nu \enbrace*{ \benbrace*{\set{F,z^{-m}},0 }_\st } = \Underbrace{\benbrace*{F_-,0}_\st}{=0}  \otimes (H-1)
	+ \benbrace*{F,0}_\st \otimes H^{-m} 
\end{align*}

\section*{Konsequenzen aus der Bott-Periodizität}

\subsection*{Definition} % (fold)
\label{sub:definition}
Sei $X$ ein kompakter Hausdorffraum mit Basispunkt $x_0$. Setze
\[
	\tilde{K}^{-n} (X) = K^{-n}(X,\pt)
\]
$\tilde{}$ heißt reduzierte $K$-Theorie von $(X,x_0)$. Theorem 3.21 impliziert
\begin{itemize}
	\item $\tilde{K}^0(S^2) \cong \mathds{Z} =  \langle \beta_2 := 1 - [H,0]_\st\rangle$ mit $\beta_2^2=0$
	\item allgemeiner $\tilde{K}^0(S^{2n}) = \langle \beta_{2m}\rangle$ mit $\beta_{2m}^2=0$, wobei $\beta_{2m}$ aus $\beta_{2m-2}$ durch äußere Multiplikation mit $\beta_2$
	entsteht:
	\[ 
		\begin{tikzcd}[]
			\tilde{K}^0(S^{2m-2}) \otimes \tilde{K}(S^2) \rar["\cong"] & K^0(D^{2m}, S^{2m-1}) \otimes K^0(D^2,S^1) \dar \\
			& K^0 \enbrace*{(D^{2m-2, S^{2m-3}}) \times (D^2,S^1)} \dar["\cong"] \\
			& K^0 \enbrace*{D^{2m}, S^{2m-1}} \cong \tilde{K}^0(S^{2m}) 
		\end{tikzcd} 
	\]
	Sei $j \colon K^0(S^2,\pt) \cong K^0(D^2,S^1) \cong K^0(I^2, \partial I^2) \cong K^{-2}(\pt)$, $\beta_2 \mapsto \hat{\beta}_2= j(\beta_2)$
\end{itemize}
% subsection definition (end)


\subsection*{Theorem} % (fold)
\label{sub:theorem}
Sei $(X,A)$ ein Ko-Raumpaar. Dann ist 
\[
	\begin{tikzcd}
		K^{-n}(X,A) \rar["\hat{\beta}_2"] & K^{-n-2}(X,A)
	\end{tikzcd}
\]
ein Isomorphismus.
\minisec{Beweis}
Es genügt $n=0$ zu betrachten. Sei $A=\emptyset$. Die exakte Sequenz
\[
	\begin{tikzcd}
		0 \rar & K(S^2,\pt) \rar & K^0(S^2) \rar & K^0(\pt) \rar & 0
	\end{tikzcd}
\]
spaltet, also spaltet auch
\[
	\begin{tikzcd}
		0 \rar & K^0(X) \otimes K^0(S^2,\pt) \rar & K^0(X) \otimes K^0(S^2) \rar & K^0(X) \otimes K^0(\pt) \rar & 0
	\end{tikzcd}
\]
Betrachte
\[
	\begin{tikzcd}
		0 \rar & K^0(X) \otimes K^0(S^2,\pt) \dar["\cong"] \rar & K^0(X) \otimes K^0(S^2) \rar \dar["\cong\text{ 3.21}",] & K^0(X) \otimes K^0(\pt) \rar \dar & 0 \\
		0 \rar & K^0(X \times (S^2,\pt)) \rar & K^0(X \times S^2) \rar & K^0(X) \rar & 0
	\end{tikzcd}
\]
Für $A \not= \emptyset$ gilt $K^{-n}$ gilt $K^{-n}(X,A) \cong \tilde{K}^{-n}(\nicefrac{X}{A})$ nach Theorem 3.9 b). Angewandt auf den Pushout
\[
	\begin{tikzcd}
		A \rar[hook] \dar & X \dar \\
		* \rar & \nicefrac{X}{A}
	\end{tikzcd}
\]
Daher genügt es den Fall $A=*$ zu betrachten. Dieser folgt analog zu dem unreduzierten Fall.

\subsection*{Definition}
Für $n>0$ und ein Ko-Raumpaar $(X,A)$ definieren wir
\[
	K^n(X,A) := K^{-n}(X,A)
\]
und die Randabbildung in der exakten Paarfolge über das nachfolgende Diagramm
\[
	\begin{tikzcd}
		K^0(A) \dar["\cong"] \rar["\partial'"] & K^1(X,A) \dar["="]\\
		K^{-2}(A) \rar["\partial"] &K^{-1}(X,A)
	\end{tikzcd}
\]
Die Paarsequenzfolge reduziert dann auf die .....-Folge
\[
	\begin{tikzcd}
		K^0(X,A) \rar & K^0(X) \rar & K^0(A) \dar\\
		K^1(A) \uar & K^1(X) \lar & K^1(X,A) \lar 
	\end{tikzcd}
\]
% subsection theorem (end)

\section*{Die $K$-Theorie von $\mathds{C}P^n$}

\subsection*{Theorem 3.2.5} % (fold)
\label{sub:theorem325}
Sei $X$ ein endlicher \CW-Komplex mit $n$ Zellen. Dann ist $K^*(X) = K^0(X) \oplus K^1(X)$ eine endlich erzeugte Gruppe mit höchstens $n$ Erzeugern. Haben alle Zellen gerade
Dimension, dann ist $K^1(X)=0$ und $K^0(X)$ ist frei abelsch mit genau einem Erzeuger für jede Zelle.
\minisec{Beweis}
Per Induktion über die Anzahl der Zellen. Für $X=\pt$ ist die Aussage richtig. Es entstehe $X$ aus $A$ durch Anheften einer $k$-Zelle für ein $k \in \mathds{N}_0$. Nach dem
Theorem 3.9 b) angewandt auf das Pushout
\[
	\begin{tikzcd}
		A \rar[hook] \dar & X \dar\\
		\pt \rar & \nicefrac{X}{A}
	\end{tikzcd}
\]
gilt $K^*(X,A) \cong K^*(\nicefrac{X}{A})$. Folglich ist $\tilde{K}^n(\nicefrac{X}{A}) \to \tilde{K}^n(X) \to \tilde{K}^n(A)$ exakt. Wir haben 
\[
	\tilde{K}^n(\nicefrac{X}{A}) \cong \tilde{K}^n(S^k) \cong K(S^k,\pt) \cong K(I^k,\partial I^k) = K^{n+k}(\pt) = \begin{cases}
		\mathds{Z}, &\text{ falls }n+k \text{ gerade}\\
		0 ,&\text{ sonst}
	\end{cases}
\]
Also ist $\tilde{K}^n(\nicefrac{X}{A})$ entweder $0$ oder $\mathds{Z}$, also erhalten wir aus der Exaktheit, dass man maximal einen Erzeuger mehr braucht, um $\tilde{K}(X)$
zu erzeugen, als man braucht um $\tilde{K}^n(A)$ zu erzeugen.

Zur zweiten Aussage: Beweis durch Induktion nach Anzahl der Zellen (wie oben). Aussage ist korrekt für $X=\pt$. Sei $A$ ein \CW-Komplex, der nur aus geraden dimensionalen 
Zellen besteht und $X$ entstehe aus $A$ durch das Anheften einer $2k$-Zelle. Wir betrachten $K^1(\nicefrac{X}{A}) \to \tilde{K}^1(X) \to \tilde{K}^1(A)$. Mit 
Induktionsannahme folgt $\tilde{K}^1(A)=0$. Weiterhin gilt $\tilde{K}^1(\nicefrac{X}{A}) \cong \tilde{K}^1(S^{2k})=0$, also folgt $\tilde{K}^1(X)=0$ und auch $K^1(X)=0$.
Damit reduziert sich die 6-Term folge auf eine kurze exakte Sequenz
\[
	\begin{tikzcd}
		0 \rar & \tilde{K}^0(\nicefrac{X}{A}) \rar & K^0(X) \rar & K^0(A) \rar & 0
	\end{tikzcd}
\]
Es gilt $\tilde{K}^0(\nicefrac{X}{A}) \cong \tilde{K}(S^{2k}) \cong \mathds{Z}$ und mit Induktionsannahme ist $k^0(A)$ eine freie abelsche Gruppe vom Rang 
$\# \text{Zellen von $A$}$. Also folgt $\tilde{K}^0(X) = \tilde{K}^0(A) \oplus \mathds{Z}$ und damit $K^0(X)\cong K^0(A) \oplus \mathds{Z}$. Also ist 
\[
	K^0(X) \cong K^0(A) \oplus \mathds{Z} \cong \mathds{Z}^{\#\text{ Zellen in $A$}} \oplus \mathds{Z} = \mathds{Z}^{\#\text{Zellen von $X$}}
\]
% subsection theorem (end)

\subsection*{Korollar} % (fold)
\label{sub:korollar}
$K^0(\mathds{C}P^n) = \mathds{Z}^{n+1}$, $K^1(\mathds{C}P^n)=0$.

\bigskip

Zur RIngstruktur auf $K^*(\mathds{C}P^n)$. Sei
\[
	L_n = \set[(v,V) \in \mathds{C}^{n+1} \otimes \mathds{C}P^n]{v \in V} 
\]
das kanonische .....Bündel über $\mathds{C}P^n$. Für $n=1$ gilt $L_1=H$ mit $H$ wie oben. Sei $\mathds{C}P^{n-1} \xhookrightarrow{p}\mathds{C}P^n$ die Inklusion. Dann gilt
$L_{n-1} \cong p^* L_n$.
% subsection korollar (end)

\subsection*{Satz} % (fold)
\label{sub:satz}
\[
	K^0(\mathds{C}P^n) \cong \nicefrac{\mathds{Z}[L]}{(L-1)^{n+1}}
\]
als Ring, induziert durch $[L_n,0]_\st \mapsfrom L$
\minisec{Beweis}
Für $n=1$ ist die Aussage richtig (siehe Satz 3.2.1). Wir betrachten $(\mathds{C}P^n, \mathds{C}P^{n-1})$ und die zugehörigen Sequenz aus der 6-Term Sequenz
\[
	\begin{tikzcd}
		0 \rar & K^0(\mathds{C}P^n, \mathds{C}P^{n-1}) \rar & K^0(\mathds{C}P^n) \rar["\rho^*"] & K^0(\mathds{C}P^{n-1}) \rar & 0
	\end{tikzcd}
\]
Wir zeigen: $(a_n)\cdot (L_n-1)^n$ erzeugt den Kern von $\rho^*$. Anmerkung: Hieraus folgt der Satz , wenn wir annehmen, dass 
$K^0(\mathds{C}P^{n-1}) \cong \nicefrac{\mathds{Z}[L]}{(L-1)^n}$, denn dann ist
\[
	\set{1, (L_n-1), (L_n-1)^2, \ldots , (L_n-1)^n} 
\]
eine additive Basis von $K^0(\mathds{C}P^n)$.

$\mathds{C}P^n \cong$ Quotient von $\nicefrac{S^{2n+1}}{S^1}$, wobei $S^1$ durch Multiplikation wirkt. Statt $S^{2n+1} \subseteq \mathds{C}^{n+1}$ können wir auch den Rand
von $D_0^2 \times D_1^2 \times \ldots \times D_n^2 \subseteq C^{n+1}$ betrachten, wobei $D_i \subseteq \mathds{C}$ die Einheitssphäre ist.
\[
	\partial \enbrace*{D_0^2 \times D_1^2 \times \ldots \times D_n^2} = \bigcup_i D_0^2 \times \ldots \times \partial D_i^2 \times D_{i+1}^2 \times \ldots \times D_n^2 
\]
Den Quotient von $D_0^2 \times \ldots \times \partial D_i^2 \times D_{i+1}^2 \times \ldots \times D_n^2 $ bezüglich der $S^1$-Operation bezeichnen wir mit 
$C_i \subseteq \mathds{C}P^n$. $C_i$ ist homöomorph zu $D_0^2 \times \ldots \times \hat{D}_i^2 \times D_{i+1}^2 \times \ldots \times D_n^2 $ und es gilt
\[
	C_i \cap C_j = \partial C_i \cap \partial C_j
\]
Betrachte $C_0 = D_1^2 \times \ldots \times D_n^2$ und $\partial C_0 = \bigcup_i \partial_i C_0$ mit 
\[
	\partial_i C_0 = D_1^1 \times \ldots \times \partial D_i^2 \times \ldots \times D_n^2 \subseteq C_i
\]
Die Inklusionen $(D_i^2,\partial D_i^2) \hookrightarrow (C_0,\partial_i C_0) \subseteq (\mathds{C}P^n,C_i)$ Wir erhalten 
\[
	\begin{tikzcd}
			K^0(D^2_0,\partial D_0^2) \otimes  \ldots \otimes K^0(D_n^2,\partial D_n^2) \rar 
			& K^0 \enbrace*{ (D_1^2,\partial D_1^2) \times \ldots \times (D_n^2, \partial D_n^2)} \\
			K^0(C_0,\partial_1 C_0) \otimes \ldots \otimes K^0(C_0,\partial_n C_0) \rar \uar["\cong"] & K^0(C_0,\partial C_0) \uar["\cong"] \\
			K^0(\mathds{C}P^n,C_1) \otimes \ldots \otimes K^0(\mathds{C}P^n,C_n) \rar \dar & K^0(\mathds{C}P^n, C_1 \cup \ldots \cup C_n) \cong 
			K^0(\mathds{C}P^n,\mathds{C}P^{n-1}) \dar \uar\\
			K(\mathds{C}P^n) \otimes \ldots \otimes K(\mathds{C}P^n) \rar & K^0(\mathds{C}P^n)
	\end{tikzcd}
\]
\todo[inline]{das brauche ich aus den Notizen…}
% subsection satz (end)
% section 3 (end)
\newpage

\section{Die Hopf-Invariante} % (fold)
\label{sec:die_hopf_invariante}
Zur Erinnerung: Wir wollen den folgenden Satz beweisen:

\Adams*
Ziel für heute: Definition der Hopf-Invariante sowie (i)$\Rightarrow $(ii).


\subsection{Definition} % (fold)
\label{sub:41}
Ein $H$-Raum ist ein punktierter topologischer Raum $(X,e)$ mit Basispunkt $e \in X$ zusammen mit einer stetigen Abbildung $\mu \colon X \times X \to X$, so dass 
\begin{enumerate}[(i)]
	\item $\mu(e,e)=e$
	\item Die Abbildung $X \to X$, $x \mapsto \mu(x,e)$ und $X \to X$, $x \mapsto \mu(e,x)$ sind beide homotop zur Identität auf $X$.
\end{enumerate}
% subsection 41 (end)

\subsection{Satz} % (fold)
\label{sub:42}
Falls $\mathds{R}^n$ die Struktur einer reellen Divisionsalgebra hat, dann existiert eine stetige Abbildung $\mu \colon S^{n-1} \times S^{n-1} \to S^{n-1}$, sodass 
$(S^{n-1},e,\mu)$ ein $H$-Raum ist.
\minisec{Beweis}
Sei $e \in \mathds{R}^n$ die Eins der Divisionsalgebra. Dann gilt $\norm{e}\not=0$. Wir können ohne Beschränkung der Allgemeinheit $\norm{e}=1 $ annehmen. Wenn dies nicht 
gilt, dann definieren wir eine neue Divisionsalgebrastruktur auf $\mathds{R}^n$ durch 
\[
	x * y := x \bullet y \cdot \norm{e} 
\]
mit Einselement $\frac{e}{\norm{e}} =: e'$. Ohne Beschränkung der Allgemeinheit ist also $e \in S^{n-1} \subset \mathds{R}^n$. Sei 
$\mu \colon S^{n-1} \times S^{n-1} \to S^{n-1}$ gegeben durch 
\[
	\mu(x,y) = \frac{x * y}{\norm{x * y} } 
\]
Dies ist wohldefiniert, da $x * y\not=0$ für $x,y \not=0$. Außerdem gilt $\mu(x,e)=x$, $\mu(e,x)=x$, \bewende
% subsection 42 (end)

\subsection{Lemma} % (fold)
\label{sub:43}
Seien $k,l >0$. Dann gilt
\[
	K^0 \enbrace*{S^{2l} \times S^{2k}} \cong K^0(S^{2l}) \otimes K^0(S^{2k}) \cong \nicefrac{\mathds{Z}[\alpha,\beta]}{(\alpha^2,\beta^2)}
\]
\minisec{Beweis}
Wir haben $\tilde{K}^0(S^{2l}) \cong \mathds{Z} = \langle \alpha \rangle$ mit $\alpha^2=0$. Die kurze exakte Sequenz 
\[
	\begin{tikzcd}
		0 \rar & K^0(S^{2l},\pt) \rar & K^0(S^{2l}) \rar & K^0(\pt)\cong \mathds{Z} \rar & 0
	\end{tikzcd}
\]
spaltet durch einen Ringhomomorphismus $\mathds{Z} \cong K^0(\pt) \to K^0(S^{2l})$. Hieraus lesen wir die Ringstruktur von $K^0(S^{2l})$ ab: 
$K^0(S^{2l})\cong \mathds{Z} + \mathds{Z} \cdot \alpha $ mit $\alpha^2=0$. Das heißt 
\[
	K^0(S^{2l}) \cong \nicefrac{\mathds{Z}[\alpha]}{(\alpha^2)}
\]
und somit $K^0{S^{2l}} \otimes K^0(S^{2k}) \cong \nicefrac{\mathds{Z}[\alpha,\beta]}{(\alpha^2,\beta^2)}$. Aus Bott-Periodizität folgt 
\begin{align*}
	K^0(S^{2k},\pt) = \tilde{K}^0(S^{2k}) \cong \tilde{K}^0(S^{2k-2}) \otimes \tilde{K}^0(S^2) \\
	K^0(X) \otimes K^0(S^2,\pt) \cong K^0(X \times (S^2,\pt))
\end{align*}
Zusammen folgt durch Induktion, dass 
\[
	K^0(X) \otimes  K^0(S^{2k},\pt ) \cong K^0 \enbrace*{X \times (S^{2k},\pt)} 
\]
Jetzt betrachten wir
\[
	\begin{tikzcd}
		0 \rar & K^0(X) \otimes  K^0(S^{2k},\pt)\dar["\cong"] \rar & K^0(X) \otimes K^0(S{2k}) \dar["(1)"] \rar & K^0(X) \otimes K^0(\pt) \dar["\cong"] \rar & 0 \\
		0 \rar & K^0 \enbrace*{X \times (S^{2k},\pt)} \rar & K^0(X \times S^{2k}) \rar & K^0(X \times \pt) \rar & 0 
	\end{tikzcd}
\]
$(1)$ ist auch ein Isomorphismus. Die Aussage folgt für $X=S^{2l}$. \bewende
% subsection 43 (end)

\subsection{Satz} % (fold)
\label{sub:44}
Falls $S^{n-1}$ ein $H$-Raum ist, dann ist entweder $n=1$ oder $n\ge 2$ gerade.
\minisec{Beweis}
Wir zeigen: Für $k>0$ kann $S^{2k}$ kein $H$-Raum sein. Angenommen wir haben $H$-Raum-Struktur $\mu \colon S^{2k} \times S^{2k} \to S^{2k}$ mit Eins $e$. 
\[
	\begin{tikzcd}
		\mu^* \colon \nicefrac{\mathds{Z}[\gamma]}{(\gamma^2)} \rar & \nicefrac{\mathds{Z}[\alpha,\beta]}{(\alpha^2,\beta^2)} \\
		K^0(S^{2k}) \rar & K^0(S^{2k} \times S^{2k})
	\end{tikzcd}
\]
\todo{$\cong$ ins Diagramm einfügen}
Jetzt betrachte
\[
	\begin{tikzcd}[row sep=0cm]
		S^{2k} \rar["i"] & S^{2k} \times S^{2k} \rar["\mu"] & S^{2k} \\
		x \rar[mapsto] & (x,e)
	\end{tikzcd}
\]
Dies ist homotop zur Identität, das heißt $i^* \circ \mu^* = \id_{K^0(S^{2k})}$. Betrachte
\[
	\begin{tikzcd}
		K^0(S^{2k} \times S^{2k}) \rar["i^*"] & K^0(S^{2k}\times \pt) \\
		K^0(S^{2k}) \otimes K^0(S^{2k}) \uar["*"] \rar & K^0(S^{2k}) \otimes K^0(\pt) \uar \\
		\nicefrac{\mathds{Z}[\alpha,\beta]}{(\alpha^2,\beta^2)} \rar["i^*"] & \nicefrac{\mathds{Z}[\gamma]}{(\gamma^2)} \\
		\alpha \rar[mapsto] & \gamma \\
		\beta \rar[mapsto] & 0
	\end{tikzcd}
\]
Seien $r,s,t \in \mathds{Z}$, so dass $\mu^*(\gamma)= r \alpha + s \beta+ t \alpha \beta$. Dann gilt $r \gamma = i^* \circ \mu^*(\gamma)=\gamma$. Also ist $r=1$. Analog folgt
$s=1$.
\[
	\Rightarrow 0 = \mu^*(\gamma^2) = \enbrace*{\mu^*(\gamma)}^2 = 2 \alpha \beta \not=0  
\]
Widerspruch! \bewende

Eine $H$-Raum-Struktur $\mu \colon S^{n-1} \times S^{n-1} \to S^{n-1}$ liefert eine stetige Abbildung $f \colon S^{2n-1} \to S^n$ definiert wir folgt: Wir haben 
$S^{2n-1} \cong \partial (D^n \times D^n) = \partial D^n \times D^n \cup D^n \times \partial D^n$ und  $S^n \cong D^n_+ \coprod D^n_n /\sim$ durch Verklebung
entlang von $S^{n-1} = \partial D^n_+ = \partial D^n_-$. $f$ ist definiert durch
\[
	\begin{tikzcd}[ampersand replacement=\&,row sep=0cm]
		f \colon \partial D^n \times D^n \cup D^n \times \partial D^n \rar \& \nicefrac{D^n_+ \coprod D^n_-}{\sim} \\
		(x,y) \rar[mapsto] \&  {\begin{cases}
			\norm{y} \cdot \mu \enbrace*{x, \frac{y}{\norm{y} } }  , &\text{ falls } (x,x) \in \partial D^n  \times D^n\\
			\norm{x} \cdot \mu \enbrace*{\frac{x}{\norm{x} },y }  , &\text{ falls } (x,y) \in D^n \times \partial D^n
		\end{cases}}
	\end{tikzcd}
\]
$f$ ist wohldefiniert und stetig. Ab jetzt nehmen wir $n=2k$ an. Sei $f \colon S^{4k-1} \to S^{2k}$ eine stetige Abbildung. $f$ ist die Anklebeabbildung einer $4k$-Zelle
$D^{4k}$ an $S^{2k}$. Sei $C_f$ durch das Pushout-Diagramm
\[
	\begin{tikzcd}
		S^{4k-1} \rar[hook] \dar["f"] & D^{4k} \dar\\
		S^{2k} \rar[hook] &C_f
	\end{tikzcd}
\]
Es gilt $\nicefrac{C_f}{S^{2k}} \cong S^{4k}$. Damit folgt
\[
	K^0(C_f,S^{2k}) \cong K^0(D^{4k}, S^{4k-1}) \cong \tilde{K}^0(S^{4k})
\]
Wegen $\tilde{K}^1(S^{4k}) \cong \tilde{K}^1(S^{2k}) =0$ erhalten wir aus der exakten Tripelsequenz für $C_f \supset S^{2k} \supset \pt$ die kurze exakte Sequenz
\[
	\begin{tikzcd}
		0 \rar & \tilde{K}^0(S^{4k}) \rar & \tilde{K}^0(C_f) \rar & \tilde{K}^0(S^{2k}) \rar & 0
	\end{tikzcd}
\]
Sei $\alpha \in \tilde{K}^0(C_f)$ das Bild des Generator von $\tilde{K}^0(S^{4k})$ in $\tilde{K}^0(C_f)$. Sei $\beta \in \tilde{K}^0(C_f)$ ein Element, welches auf den
Generator von $\tilde{K}^0(S^{2k})$ abbildet. Dann ist das Bild von $\beta^2$ in $\tilde{K}^0(S^{2k})$ gleich $0$. Also ist $\beta^2= h_f \cdot \alpha$ für ein 
$h_f \in \mathds{Z}$.
% subsection 44 (end)

\subsection{Lemma} % (fold)
\label{sub:45}
$h_f$ ist unabhängig von der Wahl von $\beta$.
\minisec{Beweis}
Sei $\beta'$ ein weiteres Element. Dann gilt $\beta' = \beta + m \cdot \alpha$ für ein $m \in \mathds{Z}$. Es gilt dann
\[
	(\beta')^2 = \beta^2 + 2m \beta \alpha
\]
Das Bild von $\alpha \beta$ in $\tilde{K}^0(S^{2k})$ ist $0$. Damit ist $\alpha \beta= r \cdot \alpha$ für ein $r \in \mathds{Z}$. Es folgt
\[
	r \alpha \beta = \alpha \beta^2 = h_f \cdot \alpha^2=0
\]
da $\alpha^2=0$ ist.  Damit ist $\alpha \beta=0$, da $\alpha \beta \in \im (\mathds{Z}= \tilde{K}^0(S^{4k}) \hookrightarrow \tilde{K}^0(C_f))$. Also gilt 
$(\beta')^2=\beta^2$. \bewende
% subsection 45 (end)

\subsection{Definition} % (fold)
\label{sub:46}
Das Element $h_f \in \mathds{Z}$ heißt die \Index{Hopf-Invariante} von $f$.
% subsection 46 (end)

\subsection{Satz} % (fold)
\label{sub:47}
Sei $f\colon S^{4k} \to S^{2k}$ die stetige Abbildung assoziiert zu einer $H$-Raum-Struktur auf $S^{2k-1}$. Dann gilt $h_f \in \set{\pm 1}$.
\minisec{Beweis}
\todo[inline]{Hier fehlt noch ganz viel}


% subsection 47 (end)

\subsection*{Die äußere Algebra}
Sei $V$ ein $n$-dimensionaler $\mathds{C}$-Vektorraum.
\[
	T(V) = \bigoplus_{n=0}^\infty V^{\otimes n}
\]
mit $V^{\otimes 0}=\mathds{C}$ ist die \Index{Tensoralgebra} von $V$ mit Multiplikation induziert vom Tensorprodukt und Einselement 
$1\in \mathds{C}=V^{\otimes 0} \subset T(V)$. $T(V)$ ist $\mathds{N}$-graduiert. Sei $I \subset T(V)$ das (graduierte) zweiseitige Ideal erzeugt von 
$v \times v \in V^{\otimes 2} \subset T(V)$ für alle $v \in V$. Das heißt in $\nicefrac{T(V)}{I}$ gilt
\begin{equation*}
	\benbrace*{v_1 \otimes \ldots \otimes v_r} = \benbrace*{v_1 \otimes  \ldots \otimes v_r + w_1 \otimes  \ldots \otimes w_k \otimes a \otimes a \otimes w_{k+3} \otimes 
	\ldots \otimes w_r} \label{eq_auss_al} \tag{\#}
\end{equation*}
Dann heißt $\Lambda^* V:= \nicefrac{T(V)}{I}$ die \Index{äußere Algebra} von $V$. Elemente in $\Lambda^*V$ schreibt man auch 
$v_1 \wedge \ldots \wedge v_r = \benbrace*{v_1 \otimes \ldots \otimes v_r}$. Aus \eqref{eq_auss_al} folgt: 
\[
	v_1 \wedge \ldots \wedge v_i \wedge v_{i+1} \wedge \ldots \wedge v_r = - v_1 \wedge \ldots \wedge v_{i+1} \wedge v_i \wedge \ldots \wedge v_r
\]
$\Lambda^* V = \bigoplus_{n=0}^\infty \Lambda^n V$ ist graduiert über die Länge der Tensoren.
\begin{beispiel}
\leavevmode
\begin{itemize}
	\item Ist $V=\mathds{C}$, so ist $T(V) = \bigoplus_{n=0}^\infty \mathds{C}$ und $\Lambda^*(V) \cong \mathds{C} \oplus \mathds{C}$.
	\item Ist $V=\mathds{C}^2$, so erhalten wir $\Lambda^* V \cong \mathds{C} \oplus \mathds{C}^2 \oplus \mathds{C}$, wobei das letze 
	$\mathds{C} = \Span \set{e_1 \wedge e_2} $ ist.
	\item Allgemeiner gilt für $V=\mathds{C}^r$
	\[
		\Lambda^* V \cong \mathds{C} \otimes C^r \oplus \ldots \oplus C^{\binom{r}{i}} \oplus \ldots \oplus \underset{\Lambda^r V}{\mathds{C}} \oplus 0 \ldots 
	\]
\end{itemize}
\end{beispiel}

Der Dualraum $(\Lambda^n V)^*$ ist kanonisch isomorph zum Vektorraum der antisymmetrischen $n$-Formen auf $V$, also Abbildungen $\varphi \colon V \times \ldots \times V \to C$, die multilinear und antisymmetrisch sind.

Die Zuordnungen $V \mapsto \Lambda^n V$ und $V \mapsto\Lambda^* V$ lassen sich zu stetigen Funktoren $\VEKT_\mathds{C} \to \VEKT_\mathds{C}$ beziehungsweise
$\VEKT_\mathds{C} \to \mathrm{Gr-}\VEKT_\mathds{C}$ erweitern. Nach \autoref{ind_funk_vektorbuendel} existieren Funktoren
\[
	\Lambda^n \colon \VEKT_\mathds{C}(X) \longrightarrow \VEKT_\mathds{C}(X) \qquad \Lambda^* \colon \VEKT_\mathds{C}(X) \longrightarrow \mathrm{Gr-}\VEKT_\mathds{C}(X)
\]
Eigenschaften dieser Funktoren: 
\begin{enumerate}[a)]
	\item $\Lambda^n(E_1 \oplus E_2) \cong \bigoplus_{i+j=n} \Lambda^i(E_1) \otimes \Lambda^j(E_2)$
	\item $\Lambda^0(E) \cong X \times \mathds{C}$, $\Lambda^1(E) \cong E$
	\item $\Lambda^n(E)=0$ für $n > \dim E$.
\end{enumerate}
\[
	\Lambda^n(E_1 \oplus E_2) = \benbrace*{\Lambda^*(E_1 \oplus E_2)}_n = \enbrace*{\Lambda^*(E_1) \otimes \Lambda^*(E_2)}_n= (*,*)  
\]

\subsection*{Elementarsymmetrische Polynome}
Sei $p \in \mathds{Z} \benbrace*{x_1, \ldots ,x_n}$ ein Polynom in Variablen $x_1, \ldots ,x_n$. Die Permutationsgruppe $\Sigma_n$ von $\set{1,\ldots ,n}$ wirkt auf $\mathds{Z}[x_1,\ldots ,x_n]$ durch
\[
	(\tau \cdot p)(x_1,\ldots ,x_n)=p \enbrace*{x_{\tau ^{-1}(1)}, \ldots , x_{\tau ^{-1}(n)}}  \quad , \quad \tau_1 \cdot (\tau_2 \cdot p) = (\tau_1 \cdot \tau_2) \cdot p
\]
Ein Polynom $p$ heißt $\Sigma_n$-invariant, falls $\tau \cdot p=p$ für alle $\tau \in \Sigma$ gilt.
\begin{beispiel}
\leavevmode
\begin{itemize}
	\item $\sigma_1(x_1,\ldots ,x_n)=x_1 +x_2 + \ldots + x_n$ ist $\Sigma_n$-invariant.
	\item $\sigma_2(x_1,x_2,x_3) = x_1 x_2 + x_1 x_3+ x_2 x_3$ ist $\Sigma_3$-invariant.
	\item $\sigma_r(x_1, \ldots ,x_n) = \sum_{1 \le i_1< \ldots < i_r \le n} x_{i_1} x_{i_2} \cdots x_{i_k}$ ist $\Sigma_n$-invariant. Diese Polynome heißen 
	\emph{elementarsymmetrische Polynome}
\end{itemize}
\end{beispiel}

Sei $q(x_1,\ldots ,x_n;t) = \prod_{i=1}^n (1+ t x_i)$. Es ist $(1+t x_1)(1+ t x_2)= 1+ t(x_1 + x_2) + t^2 x_1 x_2$. Dann ist $\sigma_r(x_1,\ldots ,x_n)$ der Koeffizient
von $q(x_1, \ldots ,x_n;t)$ vor $t^r$.
\begin{satz}
Ein Polynom $p \in \mathds{Z} \benbrace*{x_1,\ldots ,x_n}$ ist $\Sigma_n$-invariant genau dann, wenn es von der Form 
\[
	p(x_1, \ldots ,x_n) = s \enbrace*{ \sigma_1(x_1, \ldots ,x_n), \sigma_2(x_1, \ldots ,x_n) , \ldots ,  \sigma_n(x_1,\ldots ,x_n)} 
\]
für ein $s \in \mathds{Z} \benbrace*{x_1, \ldots ,x_n} $ ist.
\end{satz}

Sei $E \to X$ ein $\mathds{C}$-Vektorbündel. Betrachte 
\[
	\Lambda_t(E) := \sum_{i=0}^{n} \benbrace*{\Lambda^i(E),0}_\st  \cdot t^i = \sum_{i=0}^{n} \Lambda^i(E) t^i \in K^0(X)[t]
\]
Dann gilt $\Lambda_t(E_1 \oplus E_2) = \Lambda_t(E_1) \cdot \Lambda_t(E_2)$ wegen Eigenschaft a) von $\Lambda^n(E)$. Falls $E \cong L_1 \oplus L_2 \oplus \ldots \oplus L_n$ 
für eindimensionale $\mathds{C}$-Vektorbündel $L_i \to X$. Dann gilt
\[
	\Lambda_t(E) = \prod_{i=1}^n \Lambda_t(L_i) = \prod_{i=1}^n (1+ L_i t)
\]
Durch Koeffizientenvergleich bekommen wir 
\[
	\Lambda^i(E) \cong \sigma_i \enbrace*{L_1, \ldots , L_n} 
\]
für dieses $E$. Wir wollen $\psi^m \enbrace*{L_1 \oplus  \ldots  \oplus  L_n} = L_1^m + \ldots + L_n^m $. Das Polynom $p_m(x_1,\ldots,x_n)= x_1^m + \ldots + x_n^m$ ist 
$\Sigma_n$-invariant, das heißt es gibt $s_m(x_1,\ldots ,x_n)$ sodass $p_m= s_m(\sigma_1, \ldots ,\sigma_n)$. Jetzt definieren wir 
\[
	\psi^m(E) = s_m \enbrace*{\Lambda^1(E), \ldots, \Lambda^n(E)} 
\]
\begin{satz}[name=Spaltungsprinzip]
Sei $X$ ein kompakter Hausdorffraum. Sei $E \to X$ ein endlichdimensionales $\mathds{C}$-Vektorbündel. Dann gibt es einen kompakten Hausdorffraum $F(E)$ zusammen mit einer
stetigen Abbildung $p \colon F(E) \to X$, sodass $p^* \colon K^0(X) \to K^0 \enbrace*{F(E)}$ injektiv ist und $p^*E$ in eine SUmme von eindimensionalen 
$\mathds{C}$-Vektorbündeln zerfällt.
\end{satz}

\begin{beweis}[name=Existenz Adams Operation, 4.8]
\leavevmode
\begin{itemize}
	\item  Natürlichkeit folgt aus Funktorialität von $\Lambda^n$.
\end{itemize}
Seien $E_1, E_2$ $\mathds{C}$-Vektorraumbündel über $X$. Zu zeigen: $\psi^m(E_1 \oplus 	E_2) = \psi^m(E_1) + \psi^m(E_2)$. Sei $p_1 \colon F(E_1) \to X$ aus dem 
Spaltungsprinzip. Betrachte $p_2 \colon F(p_1^*E_2) \to F(E_1)$. Dann gilt $p_1^* E_1 \cong L^m_1 \oplus  \ldots \oplus L^m_r$ und
\[
	p_2^* p_1^* E_2 \cong L_1^{(2)} \oplus  \ldots \oplus L_s^{(2)}
\]
Es folgt \todo{hier ist noch ein Fehler im Code, deswegen auskommentiert}
% \begin{align*}
% 	p_2^* p_^* \psi^m \enbrace*{E_1 \oplus E_2} = \psi^m \enbrace*{p_2^* p_1^* \enbrace*{E_1 \oplus E_2}}
% 	&= \psi^m \enbrace*{p_2^* L_1^{(1)} \oplus \oplus p_2^* L_r^{(1)} \oplus L_1^{(2)} \oplus \ldots   \oplus L_s^{(2)}} \\
% 	&= \psi^m \enbrace*{p_2^* L_1^{(1)} \oplus  \ldots \oplus p_2^* L_r^{(1)}} + \psi^m \enbrace*{L_1^{2} \oplus  \ldots \oplus L_s^{(2)}} \\
% 	&= \psi^m \enbrace*{p_2^* p_1^* E_1} + \psi^m (p_2^* p_1^* E_2) \\
% 	&=  p_2^* p_1^* \enbrace*{\psi^m \enbrace*{E_1} + \psi^m(E_2) }
% \end{align*}
Also gilt $\psi^m \enbrace*{E_1 \oplus E_2} = \psi^m(E_1) + \psi^m(E_2)$. Jetzt definiere $\psi^m(E_1 - E_2) = \psi^m(E_1) - \psi^m(E_2)$. Die Multiplikativität geht analog.
Außerdem gilt
\begin{align*}
	\psi^m \enbrace*{\psi^l \enbrace*{L_1 \oplus L_r} } = L_1^{ml} \oplus \ldots \oplus L_r^{ml} = \psi^{ml} \enbrace*{L_1 \oplus \ldots \ldots L_r}  
\end{align*}
und
\[
	\psi^p \enbrace*{L_1 \oplus \ldots \oplus L_r} = L_1^p \oplus \ldots \oplus L_r^p = \enbrace*{L_1 \oplus \ldots \oplus L_r}^p \mod p \qedhere 
\]
\end{beweis}

\minisec{Idee des Beweises vom Spaltungsprinzip}
Betrachte das zu $E$ assoziierte projektive Bündel. Sei $s_0 \colon X \to E $ der Nullschnitt von $q \colon E \to X$.
\[
	P(E) := \nicefrac{E \setminus s_0(X)}{\sim} \qquad v \sim w \text{ falls } q(v)=q(w), v=\lambda w \text{ für ein } \lambda \in \mathds{C}^*
\]
$q' \colon P(E) \to X$, $q(w)= q'([w])$.
Sei $L= \set[{(v,[w] \in E \times P(E))}]{v=\lambda w \text{ für ein }\lambda \in \mathds{C}^*}$. Dann ist $L \to P(E)$ ein eindimensionales Vektorbündel über $P(E)$ und wir
haben eine kurze exakte Sequenz von Vektorbündeln
\[
	\begin{tikzcd}
		0 \rar & L \rar & (q')^* E \rar & Q \rar & 0
	\end{tikzcd}
\]
Es folgt nun $(q')^* E \cong L \oplus Q$. Es bleibt zu zeigen: $(q')^* \colon K^0(X) \to K^0(P(E))$ ist injektiv. Dies braucht den Thom-Isomorphismus oder das
Leray-Hirsch-Theorem.

$E=X \times \mathds{C}^{n+1}$ $\Rightarrow $ $P(E)=X \times \mathds{C}P^n$ und 
\[
	K^0(X) \otimes K^0(\mathds{C}P^n) \xrightarrow{\cong} K^0(X \times \mathds{C}P^n) 
\]

% section die_hopf_invariante (end)





\cleardoubleoddemptypage
\appendix
\section{Anhang}








\subsection{Lokale Trivialität des Tangentialbündels auf $S^n$} % (fold)
\label{sub:lokal_trivial_Sn}
Sei $S^n \subset \mathds{R}^{n+1}$ die $n$-Sphäre und sei 
\[
	\Tmap S^n = \set[(x,v) \in S^n \times \mathds{R}^{n+1}]{\skal*{x}{v}=0} 
\]
Wir definieren $p \colon \Tmap S^n \to S^n$ durch $p(x,v)=x$. Wir zeigen, dass $p \colon \Tmap S^n \to S^n$ lokal trivial ist.
\minisec{Beweis}
Sei $x_0 \in S^n$ gegeben. Sei $U := \set[x \in S^n]{\skal*{x}{x_0}>0}$. Da $\skal*{x_0}{x_0}=1$ ist, ist $x_0 \in U$. Weiter ist $U$ offen in $S^n$, also eine offene 
Umgebung von $x_0$. Sei $H \subset \mathds{R}^{n+1}$ die Hyperebene gegeben durch $H = \set[v \in \mathds{R}^{n+1}]{\skal*{x}{v}=0}$. Da dies ein $n$-dimensionaler 
Vektorraum ist, finden wir einen Isomorphismus $\kappa \colon H \to \mathds{R}^n$. Sei $\Tmap S^n \big|_U=p^{-1}(U)$. Wegen $\skal*{x_0}{x_0}=1$ liegt $v- \skal*{x_0}{v} x_0$
in $H$ für alle $v \in \mathds{R}^{n+1}$. Also können wir definieren:
\[
	\varphi \colon \Tmap S^n\big|_{U} \to U \times \mathds{R}^n \quad , \quad (x,v) \mapsto \enbrace[\big]{x, \kappa(v - \skal*{x_0}{v} x_0)} 
\]
Außerdem definieren wir $\psi \colon U \times \mathds{R}^n \to \Tmap S^n \big|_U$ durch
\[
	\psi(x,w) = \enbrace*{x , \kappa^{-1}(w) - \frac{\skal*{x}{\kappa^{-1}(w)}}{\skal*{x}{x_0}} x_0 } 
\]
Dies ist wohldefiniert, da für $x \in U$ der Nenner stets ungleich $0$ ist und weiter
\[
	\skal*{x}{\kappa^{-1}(w) - \frac{\skal*{x}{\kappa ^{-1}(w)}}{\skal*{x}{x_0}} x_0} = \skal*{x}{\kappa^{-1}(w)} - \frac{\skal*{x}{\kappa ^{-1}(w)}}{\skal*{x}{x_0}} 
	\skal*{x}{x_0} = 0
\]
gilt. Weiter sind beide Abbildungen offensichtlich stetig. Für die Komposition gilt
\begin{align*}
	(\psi \circ \varphi)(x,v) = \enbrace*{x, v - \skal*{x_0}{v} x_0 - \frac{\skal*{x}{v- \skal*{x_0}{v} x_0}}{\skal*{x}{x_0}}x_0 }
	= \enbrace*{x, v - \skal*{x_0}{v} x_0 + \skal*{x_0}{v}x_0} = (x,v)  
\end{align*}
Dabei haben wir benutzt, dass $\skal*{x}{v}=0$ gilt. Umgekehrt gilt
\begin{align*}
	(\varphi \circ \psi)(x,w) &= \enbrace*{x , \kappa \enbrace*{\kappa^{-1}(w) - \frac{\skal*{x}{\kappa^{-1}(w)}}{\skal*{x}{x_0}}x_0 - \skal*{x_0}{\kappa^{-1}(w) - 
	\frac{\skal*{x}{\kappa^{-1}(w)}}{\skal*{x}{x_0}}x_0} x_0 } } \\
	&= \enbrace*{x , \kappa \enbrace*{\kappa^{-1}(w) - \frac{\skal*{x}{\kappa^{-1}(w)}}{\skal*{x}{x_0}}x_0 - 
	\skal*{x_0}{\kappa^{-1}(w)} x_0 + \frac{\skal*{x}{\kappa^{-1}(w)}}{\skal*{x}{x_0}} x_0 }}\\
	&= \enbrace[\Big]{x , \kappa \enbrace*{\kappa^{-1}(w) - \skal*{x_0}{\kappa^{-1}(w)} x_0} } = (x,w)
\end{align*}
da wieder $\skal*{x_0}{\kappa ^{-1}(w)}=0$, weil $\kappa ^{-1}(w) \in H$ ist. Damit ist $\varphi$ ein Homöomorphismus. Es ist außerdem klar, dass $\varphi$ eingeschränkt auf 
die Fasern eine lineare Abbildung ist. Damit ist das Tangentialbündel $\Tmap S^n$ lokal trivial. \bewende
% subsection lokal_trivial_Sn (end)








\cleardoubleoddemptypage
\pagenumbering{Alph}
\setcounter{page}{1}

\printindex
\listoffigures
\todototoc
\listoftodos[To-do's und andere Baustellen]
\end{document}
